\documentclass[11pt,a4paper]{article}
\usepackage[margin=2.3cm]{geometry}
\usepackage[T1]{fontenc}
\usepackage{lmodern}
\usepackage{amsmath,amssymb,amsthm,mathtools}
\usepackage{bm}
\usepackage{booktabs}
\usepackage{array}
\usepackage{longtable}
\usepackage{enumitem}
\usepackage{xcolor}
\usepackage[hidelinks]{hyperref}
\usepackage{microtype}
\usepackage{xurl}
\emergencystretch=3em

\numberwithin{equation}{section}
\newtheorem{theorem}{Theorem}[section]
\newtheorem{proposition}[theorem]{Proposition}
\newtheorem{lemma}[theorem]{Lemma}
\newtheorem{corollary}[theorem]{Corollary}
\theoremstyle{definition}
\newtheorem{definition}[theorem]{Definition}
\newtheorem{assumption}[theorem]{Assumption}
\newtheorem{example}[theorem]{Example}
\theoremstyle{remark}
\newtheorem{remark}[theorem]{Remark}

\newcommand{\E}{\mathbb{E}}
\newcommand{\Var}{\operatorname{Var}}
\newcommand{\Cov}{\operatorname{Cov}}
\newcommand{\R}{\mathbb{R}}
\newcommand{\N}{\mathbb{N}}
\newcommand{\PP}{\mathbb{P}}
\newcommand{\ind}{\mathbf{1}}
\newcommand{\CRPS}{\operatorname{CRPS}}
\newcommand{\VS}{\operatorname{VS}}
\newcommand{\PB}{\operatorname{PB}}
\newcommand{\Snorm}{S^{\mathrm{norm}}}
\newcommand{\MO}{\mathrm{M0}}
\newcommand{\sd}{\operatorname{sd}}
\newcommand{\diag}{\operatorname{diag}}
\newcommand{\sgn}{\operatorname{sgn}}
\newcommand{\argmin}{\operatorname*{arg\,min}}
\newcommand{\argmax}{\operatorname*{arg\,max}}
\newcommand{\code}[1]{\texttt{#1}}

\title{\textbf{A Ratio-Aware, Gated Probabilistic Forecasting Engine\\for Reasoning-Augmented Time-Series Forecasting}\\[4pt]
\large Agenthon 2026, Track 2 --- Technical Report and Mathematical Specification\\
\normalsize Implementation \code{agenthon-track2} v1.0.1}
\author{Team 89}
\date{29 September 2026}

\begin{document}
\maketitle

\begin{abstract}
\noindent
We describe, specify and validate a complete forecasting system for Agenthon~2026 Track~2, in which an agent reads point-in-time
numerical panels and a dated text corpus and must emit $N$ joint Monte-Carlo draws over a grid of (asset, horizon) cells.
Submissions are ranked by the equal-weight mean over cards of a composite of marginal CRPS, a variogram score and a pinball tail
penalty, each normalised \emph{per card} by the same component of an official text-blind Gaussian random-walk baseline (M0),
clipped at $4$. We show that this normalisation turns forecasting into a decision problem under a \emph{ratio loss}, and derive
three consequences that dominate the design: (i) a fully reproducible replica of M0 is simultaneously the optimal fail-safe
(it scores exactly $1$) and the reference against which hypothetical outcomes can be scored; (ii) for any proper score, the
metric-optimal submission is the honest predictive belief reweighted by the inverse of the baseline's score, which we implement as
a tempered importance-resampling operator with an honesty floor; (iii) because CRPS and the variogram score are convex in the
forecast distribution, a mixture with the baseline's own law caps the loss on cards where the baseline happens to be right.
Around this decision layer we build a strong but humble numerical prior (a mean-reverting variance term structure with a shrunk
variance ratio, path-coherent joint sampling and a soft lower bound), a lexical, per-unit text pipeline feeding one
schema-constrained call to the organizer-hosted language model, and an information gate that provably returns the numerical prior
when the text carries no information. A contextual policy with pessimistic safe-policy-improvement chooses the decision action per
card. Every fitted quantity is \emph{vintaged} by the last date of data it saw, so that the whole 2003--2024 walk-forward evaluation
is out of sample. On 3{,}896 walk-forward cards the shipped system scores $0.968$ (95\% block-bootstrap interval $[0.952,0.985]$);
the first official Development submission scored $0.9782$ on $71/71$ cards. We report every ablation, including the components that
were rejected, and state plainly what could not be measured.
\end{abstract}

\tableofcontents
\newpage

\section{The forecasting problem}\label{sec:problem}

\subsection{Cards, information sets and the object to be produced}

A \emph{card} (unit) is a tuple
\begin{equation}
  C = \bigl(\tau,\ \mathcal{A},\ \mathcal{H},\ \text{type},\ \text{family},\ \bm w,\ \mathcal{L}\bigr),
\end{equation}
where $\tau$ is the as-of date (the trusted cutoff \code{[provenance] data\_cutoff}), $\mathcal{A}=\{a_1,\dots,a_d\}$ the target
assets, $\mathcal{H}$ the declared horizon keys, $\text{type}\in\{\text{level},\text{log\_return}\}$, $\text{family}\in\{\text{F1},\dots,\text{F4}\}$,
$\bm w=(w_m,w_j,w_t)=(0.5,0.3,0.2)$ the component weights and $\mathcal{L}=\{0.01,0.05,0.95,0.99\}$ the tail levels.
The \emph{grid} is the set of cells $\mathcal{G}=\mathcal{A}\times\mathcal{H}$, of size $D=|\mathcal{A}|\,|\mathcal{H}|$. Throughout, cells
are enumerated in \emph{M0 order}: sorted by asset identifier, then by horizon ascending. This is the order in which the official
baseline assigns standard-normal columns to cells, and keeping it everywhere makes our draw matrices align with the baseline's.

The agent receives two information sets, both truncated by the organizers' staging gates:
\begin{align}
  \mathcal{X}_{\le\tau} &= \{(t,a,x_{t,a}) : t\le\tau\} &&\text{(numerical panels)},\\
  \mathcal{T}_{\le\tau} &= \{(t^{\mathrm{avail}}_k, \text{doc}_k) : t^{\mathrm{avail}}_k \le \tau\} &&\text{(dated corpus)}.
\end{align}
It must output $N$ joint draws $\bm Y^{(1)},\dots,\bm Y^{(N)}\in\R^D$, $200\le N\le 20000$, where a draw index denotes one coherent
joint scenario across all cells. Formally the target of estimation is the conditional law
\begin{equation}\label{eq:target}
  P_1 \;=\; \mathcal{L}\bigl(\bm Y_{\tau+\mathcal{H}} \,\big|\, \mathcal{X}_{\le\tau},\mathcal{T}_{\le\tau}\bigr),
  \qquad
  P_0 \;=\; \mathcal{L}\bigl(\bm Y_{\tau+\mathcal{H}} \,\big|\, \mathcal{X}_{\le\tau}\bigr),
\end{equation}
where $P_0$ is the text-blind prior. The \emph{text uplift} on a card is $\Delta_{\mathrm{text}}(C)=S(P_0)-S(P_1)$, positive when the text
helped.

\subsection{Target semantics}

For a \emph{level} target the value at horizon $h$ is the panel value $h$ steps after $\tau$ (business days for daily panels, calendar
months for the four monthly macro cards; see Section~\ref{sec:baseline}). For a \emph{log\_return} target the panel rows are decimal
per-step \emph{simple} returns $r_t$ and the target is the cumulative log return
\begin{equation}
  Y_{\tau,h} \;=\; \sum_{i=1}^{h}\log(1+r_{\tau+i}),
\end{equation}
anchored at zero: the last observed return belongs to the history, not to the future total.

\subsection{Families}

The four families differ in the role of text and in the score component that discriminates:
F1 (continuation with context; marginal CRPS), F2 (text-cued regime shift, including transfer to emerging-market currencies whose
path is withheld; CRPS and tail), F3 (cross-asset reasoning; the variogram term), F4 (tail and shock from text; the pinball term).
The public practice roster has $104$ directories ($103$ cards plus one exemplar): $60$ single-cell cards, horizons concentrated on
$\{21,63,126,189\}$, and $88$ level versus $16$ log-return targets. A structural fact that matters later is that practice cards are
\emph{authored around events}: the as-of date sits just before a notable episode. We treat this as evidence that the sealed evaluation
cards (H2~2025--Q2~2026) are event-selected in the same way.

\subsection{Constraints that shape the architecture}

\begin{enumerate}[leftmargin=*]
  \item \textbf{Network.} At scoring time the only egress is an audited proxy to one organizer-hosted open model
  (\path{nvidia/nemotron-3-super-120b-a12b}, snapshot \code{rl-030326-fp8}) at \path{$MODEL_ENDPOINT/v1/chat/completions}, with
  $25$ admitted requests per card and at most $4{,}000$ output tokens per request. Thinking is on by default and consumes that budget.
  \item \textbf{Artifacts.} Pretrained neural checkpoints (time-series foundation models, embedding models, rerankers) require separate
  organizer approval; static dictionaries and BM25-type indexes are permitted. Every fitted quantity must have been computable from data
  available by each relevant task cutoff.
  \item \textbf{Runtime.} A read-only root filesystem, a non-root user with no home, a $64$\,MiB non-executable \code{/tmp}, $256$ processes,
  and a phase-shared clock of $43{,}200$\,s for about one hundred cards.
  \item \textbf{Output.} Exactly three files in \code{/output}: \code{forecast.parquet} with columns
  \code{draw:int32}, \code{asset:string}, \code{horizon:int32}, \code{value:float64}, \code{forecast\_meta.json} whose \code{unit\_id} and \code{asof} must equal the
  card's, and a non-blank \code{forecast\_rationale.md} that is never scored. Four admissibility gates g0--g3 must pass; a failure is scored
  as the worst value, $4$.
\end{enumerate}

\subsection{Summary of the approach}

The system is a directed acyclic pipeline, L0 through L23, wrapped in a fail-safe:
\begin{align*}
&\underbrace{C}_{\text{L0}}\to\underbrace{(\mathcal{X}_{\le\tau},\mathcal{T}_{\le\tau})}_{\text{L1}}\to
\underbrace{s_\tau,\pi_\tau}_{\text{L2--L4}}\to\underbrace{P_0}_{\text{L5--L7}}\to
\underbrace{E\to\bar E\to q\to R\to Z}_{\text{L8--L13}}\\
&\qquad\to\underbrace{z,\kappa,g}_{\text{L14--L16}}\to
\underbrace{P_1}_{\text{L17--L19}}\to\underbrace{\{\bm Y^{(k)}\}}_{\text{L21}}\to\underbrace{Q}_{\text{L22b}}\to\text{L23}.
\end{align*}
The distinctive element is L22b, the \emph{score-aware decision layer}, which Section~\ref{sec:ratio} derives from the scoring rule.

\section{The scoring rule, exactly}\label{sec:scoring}

\subsection{Components}

Let $\hat F$ denote the empirical law of the $N$ submitted draws $\{\bm x^{(k)}\}$ and $\bm y\in\R^D$ the realised vector.

\paragraph{Marginal CRPS.} For one cell with draws $x_1,\dots,x_N$,
\begin{equation}\label{eq:crps}
  \CRPS(\hat F,y) \;=\; \frac1N\sum_{k}|x_k-y| \;-\; \frac{1}{2N(N-1)}\sum_{k,l}|x_k-x_l|,
\end{equation}
the \emph{fair} (almost unbiased) ensemble estimator; the marginal component is the mean over the $D$ cells.

\begin{lemma}[Sorted-sum identity]\label{lem:sorted}
If $x_{(1)}\le\dots\le x_{(N)}$ then $\sum_{k,l}|x_k-x_l|=2\sum_{i=1}^N(2i-N-1)\,x_{(i)}$.
\end{lemma}
\begin{proof}
Each unordered pair $i<j$ contributes $x_{(j)}-x_{(i)}$ twice. In $\sum_{i<j}(x_{(j)}-x_{(i)})$ the element $x_{(i)}$ appears with sign $+$ in
$i-1$ pairs and with sign $-$ in $N-i$ pairs, so its coefficient is $(i-1)-(N-i)=2i-N-1$.
\end{proof}
The identity gives $O(N\log N)$ evaluation. For $K$ hypothetical outcomes at once we also need the first term fast: with prefix sums
$c_j=\sum_{i\le j}x_{(i)}$ and $k=\#\{i:x_{(i)}<y\}$,
\begin{equation}
  \sum_i|x_i-y| \;=\; (yk-c_k) + \bigl(c_N-c_k-y(N-k)\bigr),
\end{equation}
so each outcome costs a binary search. This is how the decision layer scores thousands of candidate outcomes per card.

\paragraph{Variogram score of order $p=\tfrac12$.} With unit weights over \emph{all ordered} pairs of cells,
\begin{equation}\label{eq:vs}
  \VS(\hat F,\bm y) \;=\; \sum_{i,j=1}^{D}\Bigl(|y_i-y_j|^{p}-\E_{\hat F}|X_i-X_j|^{p}\Bigr)^2 .
\end{equation}
Pairs include the same asset at two horizons and different assets at different horizons, so both cross-asset and cross-horizon geometry
are scored.

\begin{proposition}[What the variogram rewards]\label{prop:vs-opt}
For fixed truth $\bm Y\sim P$, $\E_P\VS(F,\bm Y)$ is minimised over forecasts $F$ exactly when
$\E_F|X_i-X_j|^p=\E_P|Y_i-Y_j|^p$ for every pair; the minimum equals $\sum_{i,j}\Var_P(|Y_i-Y_j|^p)$.
\end{proposition}
\begin{proof}
For each pair, $\E_P(U-c)^2$ with $U=|Y_i-Y_j|^p$ and constant $c=\E_F|X_i-X_j|^p$ equals $\Var U+(\E U-c)^2$, minimised at $c=\E U$.
The pairs decouple.
\end{proof}
Hence VS is proper but not strictly proper: it scores only the $p$-th moments of pairwise distances. Two consequences follow. First, a
common location shift of all cells does not change VS, which makes it complementary to CRPS. Second, independent sampling across
assets inflates $\E|X_i-X_j|^p$ for positively dependent pairs, and is penalised.

\paragraph{Tail pinball penalty.} With $\hat q_\alpha$ the empirical (linear-interpolation) $\alpha$-quantile of a cell,
\begin{equation}\label{eq:pb}
  \PB(\hat F,\bm y) \;=\; \frac{1}{|\mathcal L|\,D}\sum_{\alpha\in\mathcal L}\sum_{i=1}^{D}\rho_\alpha\bigl(y_i-\hat q_{\alpha,i}\bigr),\qquad
  \rho_\alpha(u)=u\bigl(\alpha-\ind\{u<0\}\bigr).
\end{equation}

\subsection{Composite and per-card normalisation}

Write $\bm S(F,\bm y)=(S_m,S_j,S_t)$ for the three raw components and $\bm S^{0}(\bm y)$ for those of the baseline M0 on the same card and
outcome. The leaderboard statistic of a card is
\begin{equation}\label{eq:snorm}
  \Snorm(F,\bm y) \;=\; \min\!\Biggl\{4,\ \sum_{k\in\{m,j,t\}} \tilde w_k\,\frac{S_k(F,\bm y)}{\bar S^{0}_k(\bm y)}\Biggr\},
  \qquad \bar S^0_k = \begin{cases} S^0_k & S^0_k>0\\ 1 & S^0_k\le 0,\end{cases}
\end{equation}
where on a single-cell card the variogram is identically $0$ and the weights are renormalised to
$\tilde{\bm w}=(0.5/0.7,\,0,\,0.2/0.7)=(0.714,0,0.286)$; otherwise $\tilde{\bm w}=\bm w$. The ranking is the equal-weight mean of
$\Snorm$ over cards, lower is better, and an inadmissible, erroring or unattempted card counts as $4$.

\begin{remark}
Three properties of \eqref{eq:snorm} drive everything that follows.
(a) The denominator is a random variable: the baseline's error \emph{on this particular outcome}.
(b) The statistic is bounded below by $0$ and above by $4$, but the ratio of a good forecast to M0 is rarely below about $0.3$, while it can
reach $4$ on cards where M0 is sharp and right. The per-card loss is therefore strongly asymmetric.
(c) The mean over cards of a ratio is not the ratio of means; heterogeneous cards are averaged on one scale because each is normalised
by its own baseline.
\end{remark}

\subsection{Implementation parity}

The system reimplements all three components in vectorised form (\code{agenthon.decision.m0\_scoring}) so that one forecast can be scored
against $K$ outcomes at once. Parity against the official \code{qfbench2\_common} and track modules was checked on random inputs:
\begin{center}
\begin{tabular}{lc}
\toprule
component & $\max|\text{ours}-\text{official}|$\\
\midrule
marginal CRPS & $3.3\times10^{-16}$\\
variogram ($p=\tfrac12$) & $0$\\
pinball (mean over levels and cells) & $0$\\
\bottomrule
\end{tabular}
\end{center}

\section{The official baseline M0 and its exact replica}\label{sec:baseline}

\subsection{Specification}

M0 is a text-blind joint Gaussian random walk, published in full except for its code and its per-card scales. For each target asset $a$:

\begin{enumerate}[leftmargin=*]
\item \textbf{Window.} Take the rows at or before $\tau$ from the first panel file (sorted filename order) containing $a$, and keep the
last $300$ observations $v_{a,1},\dots,v_{a,n}$ with dates $t_1<\dots<t_n$.
\item \textbf{Steps.} For a level target $\Delta_{a,i}=v_{a,i}-v_{a,i-1}$; for a log-return target the row itself, $\Delta_{a,i}=v_{a,i}$
(already a per-step return), with a tripwire refusing panels whose median absolute value is at least $0.2$. A step is dated at the row it ends on.
\item \textbf{Gap rule.} Drop any step spanning an interval longer than
\begin{equation}
  \ell \;=\; \max\bigl(10\cdot\mathrm{median}_i(t_i-t_{i-1}),\ 5\ \text{days}\bigr).
\end{equation}
On transfer cards this stops a withheld decade from being read as one step.
\item \textbf{Alignment.} Intersect the step dates across assets; drop a date for all assets if any asset lacks a step there.
\item \textbf{Moments.} $\mu_a=\overline{\Delta_a}$, $\Sigma=\widehat{\Cov}(\Delta)$ with the sample convention, over the intersected rows.
\item \textbf{Anchor.} $v_{a,n}$ for a level target (the last observation, not the as-of), $0$ for a log-return target.
\item \textbf{Horizon to steps.} For cell $(a,h)$ with a known target date $T$: estimate the spacing $\bar\delta$ as the mean non-hole interval;
count $s=12(y_T-y_n)+(m_T-m_n)$ calendar months if $\bar\delta>20$ days, else $s=\mathrm{round}((T-t_n)/\bar\delta)$; if $s\le 0$ or the target
date is unknown use $s=h$; otherwise keep $h$ unless $\max(s,h)/\max(\min(s,h),1)\ge 2$.
\item \textbf{Mean and covariance over cells.} With cells $c=(a_c,h_c)$ and step counts $s_c$,
\begin{equation}\label{eq:m0cov}
  m_c = \text{anchor}_{a_c} + s_c\,\mu_{a_c}, \qquad K_{cc'} = \min(s_c,s_{c'})\,\Sigma_{a_c a_{c'}} .
\end{equation}
\item \textbf{Draws.} $K\leftarrow K+10^{-10}I$; $L=\mathrm{chol}(K+10^{-9}I)$, falling back to $\diag(K+10^{-9}I)^{1/2}$;
seed $=\mathrm{crc32}(\text{unit\_id})\wedge(2^{31}-1)$; $Z\sim\mathcal N(0,I)^{500\times D}$ from \code{numpy.default\_rng(seed)}; draws
$m+ZL^\top$.
\end{enumerate}

\begin{proposition}[Path structure]\label{prop:m0path}
Under \eqref{eq:m0cov}, for one asset observed at $s_1<s_2$ steps, the increment between the two cells is independent of the first cell and
has variance $(s_2-s_1)\Sigma_{aa}$; the correlation of the two cells is $\sqrt{s_1/s_2}$.
\end{proposition}
\begin{proof}
$\Cov(X_1,X_2-X_1)=\min(s_1,s_2)\Sigma-\min(s_1,s_1)\Sigma=0$ and
$\Var(X_2-X_1)=s_2\Sigma-2s_1\Sigma+s_1\Sigma=(s_2-s_1)\Sigma$; jointly Gaussian and uncorrelated implies independent.
Finally $\mathrm{Corr}=s_1\Sigma/\sqrt{s_1s_2}\Sigma=\sqrt{s_1/s_2}$.
\end{proof}

\subsection{Verification of the replica}

The replica (\code{numerical.random\_walk.m0\_forecast}) reproduces every figure of the published worked example: the exemplar's seed is
$795{,}546{,}941$; its trailing window has $129$ rows per asset, hence $128$ steps; the counted step number is $24$ against a declared
$21$, ratio $1.14<2$, so $21$ is kept. On the monthly macro cards the month-counting branch fires and gives $8$ and $9$ steps for
declared horizons $140$ and $160$, and $2$ steps for the NFP and CPI-vintage cards. A development bug is worth recording because it would
have been silent: under \code{pandas}~3 datetime indices are not guaranteed nanosecond resolution, so interval computations of the form
$\text{asi8}/86{,}400\times10^9$ returned counts in the tens of thousands; all day arithmetic now goes through a resolution-independent
helper.

\begin{proposition}[The replica scores exactly one]\label{prop:m0one}
If the submitted draws are bit-identical to M0's, then $\Snorm=\sum_k\tilde w_k=1$ for every outcome $\bm y$.
\end{proposition}
\begin{proof}
Each ratio in \eqref{eq:snorm} is $S_k/S_k=1$ when $S_k^0>0$, and $S_k/1=S_k$ with $S_k=S^0_k\le0$ when $S^0_k\le0$. The latter happens only for the
variogram on a single-cell card, where it is $0$ and its weight is $0$. Hence the sum is $\sum_k\tilde w_k=1<4$.
\end{proof}
In the walk-forward (Section~\ref{sec:validation}) submitting the replica scores $1.0000$ on all $3{,}896$ cards, an end-to-end check of the
normalisation chain.

\subsection{Why M0 is both the fail-safe and the reference}

A card that fails costs $4$; the reference implementation shipped by the organizers averages about $1.29$. Submitting M0's own draws costs
exactly $1$ by Proposition~\ref{prop:m0one}, so the replica is the best possible fail-safe that needs no model. The same object is the
denominator of the loss: since the draws are reproducible from the card alone, the agent can compute $\bm S^0(\bm y)$ for any hypothetical
$\bm y$. Section~\ref{sec:ratio} exploits this.

\subsection{Known weaknesses of M0}

Relative to a careful forecaster M0 (i) uses a raw $300$-observation mean drift, mostly estimation noise for near-martingale series;
(ii) uses one sample variance with no conditional dynamics; (iii) is Gaussian; (iv) has no lower bound for yields; (v) on transfer cards
estimates its scale from a decades-old window. Section~\ref{sec:results} shows that correcting (i)--(iii) in the \emph{honest} forecast does not by
itself beat M0 under \eqref{eq:snorm}. The reason is the subject of the next section.

\section{Decision theory under a ratio loss}\label{sec:ratio}

\subsection{Setting}

Fix a card. Let $P$ be the agent's honest predictive belief for $\bm Y$, $F^0$ the (known) law of M0's draws, and $Q$ the law of the submitted
draws. The expected leaderboard loss of submitting $Q$ is
\begin{equation}\label{eq:risk}
  \mathcal R(Q) \;=\; \E_{\bm Y\sim P}\Bigl[\min\Bigl\{4,\ \sum_k \tilde w_k\, \frac{S_k(Q,\bm Y)}{\bar S_k^0(\bm Y)}\Bigr\}\Bigr].
\end{equation}
Ignoring the clip for the moment, $\mathcal R(Q)=\sum_k\tilde w_k\,\E_P[\omega_k(\bm Y)S_k(Q,\bm Y)]$ with outcome weights
$\omega_k(\bm y)=1/\bar S^0_k(\bm y)$ that do not depend on $Q$.

\subsection{The Bayes act is a reweighted belief}

\begin{definition}[Weighted belief]
For a measurable $\omega\ge0$ with $0<Z_\omega=\E_P\omega(\bm Y)<\infty$, the $\omega$-tilted belief is $P_\omega(dy)=\omega(y)P(dy)/Z_\omega$.
\end{definition}

\begin{theorem}[Bayes act for one weighted proper score]\label{thm:bayes}
Let $S$ be a proper scoring rule and $\omega$ as above. Then
\begin{equation}
  \argmin_Q \E_P[\omega(\bm Y) S(Q,\bm Y)] \;\ni\; P_\omega,
\end{equation}
and if $S$ is strictly proper the minimiser is unique.
\end{theorem}
\begin{proof}
For any $Q$, $\E_P[\omega S(Q,\cdot)]=\int S(Q,y)\,\omega(y)P(dy)=Z_\omega\,\E_{P_\omega}[S(Q,\bm Y)]$. Since $Z_\omega>0$ does not depend on $Q$,
minimising the left side is minimising the expected score under $P_\omega$, whose minimiser is $P_\omega$ by propriety (uniquely under strict
propriety).
\end{proof}

\begin{corollary}
CRPS is strictly proper on distributions with finite first moment, so if the composite were the marginal CRPS alone, the metric-optimal
submission would be $P$ tilted by $1/\CRPS(F^0,\bm y)$. The pinball loss at level $\alpha$ is proper for the $\alpha$-quantile functional, so the
optimal tail quantiles are the quantiles of $P$ tilted by $1/\PB(F^0,\bm y)$. The variogram score (Proposition~\ref{prop:vs-opt}) is optimised by
matching $\E|X_i-X_j|^p$ under $P$ tilted by $1/\VS(F^0,\bm y)$.
\end{corollary}

The three components call for three different tilts, and no single $Q$ attains all three optima simultaneously. We therefore restrict to a
one-parameter family and optimise that parameter empirically.

\begin{definition}[Tempered score-aware operator]\label{def:tempered}
Let $\bar s_k=\E_{P}[S^0_k(\bm Y)]$ be the average baseline component under the belief. Define the combined proxy weight and the tempered
tilt, for $\gamma\ge0$:
\begin{equation}\label{eq:weight}
  \omega(\bm y) \;=\; \Bigl(\sum_k \tilde w_k\,\frac{S^0_k(\bm y)}{\bar s_k}\Bigr)^{-1},\qquad
  Q_\gamma(d\bm y)\;\propto\;\omega(\bm y)^{\gamma}\,P(d\bm y).
\end{equation}
\end{definition}
$\gamma=0$ is the honest belief; $\gamma=1$ is the Bayes act of Theorem~\ref{thm:bayes} for the combined proxy weight. Tempering ($\gamma<1$)
hedges against the fact that $P$ is itself estimated; $\gamma>1$ over-weights the region where M0 is sharp.

\subsection{Why the tilt concentrates the submission}

\begin{proposition}[Shape of the weight for a Gaussian baseline]\label{prop:shape}
Let one cell have baseline $F^0=\mathcal N(m,\sigma^2)$ and $z=(y-m)/\sigma$. Then
$\CRPS(F^0,y)=\sigma\bigl[z(2\Phi(z)-1)+2\varphi(z)-\pi^{-1/2}\bigr]$, which is even in $z$, increasing in $|z|$, equal to
$\sigma(\sqrt2-1)/\sqrt\pi\approx0.234\sigma$ at $z=0$ and asymptotic to $\sigma|z|$. Consequently $\omega\propto 1/\CRPS$ is maximal at the baseline
median and decays like $1/|z|$.
\end{proposition}
\begin{proof}
The closed form is the standard Gaussian CRPS (Gneiting and Raftery). Its derivative in $z$ is $\sigma(2\Phi(z)-1)$, which has the sign of $z$;
evenness follows from $\Phi(-z)=1-\Phi(z)$ and the evenness of $\varphi$. At $z=0$ the bracket is $2\varphi(0)-\pi^{-1/2}=(\sqrt2-1)/\sqrt\pi$.
As $|z|\to\infty$, $2\Phi(z)-1\to\sgn z$ and $\varphi\to0$.
\end{proof}
Multiplying $P$ by a weight decaying like $|z|^{-\gamma}$ both sharpens it and pulls it toward M0's centre. The composite ratio penalises
deviations from M0's centre convexly (being wrong where M0 is right costs up to $4$; being right where M0 is wrong gains at most about
$0.7$), and the tilt is the formal counterweight to that convexity.

\subsection{A simulation with a correctly specified baseline}

Take $Y\sim\mathcal N(0,1)$, $F^0=\mathcal N(0,1)$ from $500$ draws, single-cell weights $(0.714,0.286)$, and an honest belief equal to the truth.
Averaged over $300$ replications of $4{,}000$ outcomes:
\begin{center}
\begin{tabular}{lcccc}
\toprule
& $\gamma=0$ (honest) & $\gamma=0.5$ & $\gamma=1$ & $\gamma=1.5$\\
\midrule
Gaussian truth & $1.004$ & $\mathbf{0.962}$ & $0.975$ & $1.029$\\
Student-$t_5$ truth (belief $=$ truth) & $0.976$ & $\mathbf{0.929}$ & $0.942$ & $0.992$\\
\bottomrule
\end{tabular}
\end{center}
Even when M0 is exactly right, the honest forecast scores about $1$, while the tempered tilt scores about $0.96$. The optimum is interior because
the tail term and the $1/x$ convexity pull in different directions. In real data the optimum sits higher ($\gamma\approx1$--$1.25$), because the
belief and the baseline differ (Section~\ref{sec:results}).

\subsection{Model averaging with the baseline}

\begin{lemma}[Convexity in the forecast]\label{lem:convex}
For fixed $y$, $F\mapsto\CRPS(F,y)$ is convex on distributions with finite mean, and so is $F\mapsto\VS(F,\bm y)$.
\end{lemma}
\begin{proof}
$\CRPS(F,y)=\E_F|X-y|-\frac12\E_{F\otimes F}|X-X'|$. The first term is linear in $F$. For $F_\lambda=(1-\lambda)F+\lambda G$, the second is
$\frac12 B(F_\lambda,F_\lambda)$ with the symmetric bilinear form $B(F,G)=\iint|x-x'|F(dx)G(dx')$. For signed measures of total mass zero,
$B(\nu,\nu)=-2\int(\nu(-\infty,t])^2\,dt\le0$ (the energy-distance identity), so $B$ is negative semidefinite on differences, and
$\lambda\mapsto -\frac12B(F_\lambda,F_\lambda)$ is convex. For VS, $\E_{F}|X_i-X_j|^p$ is linear in $F$ and $c\mapsto(u-c)^2$ is convex, so each summand
is a convex function of a linear map of $F$.
\end{proof}

\begin{theorem}[Mixing caps the ratio]\label{thm:mix}
Let $Q_\lambda=(1-\lambda)Q+\lambda F^0$, $\lambda\in[0,1]$. For every outcome $\bm y$ and $k\in\{m,j\}$,
\begin{equation}
  \frac{S_k(Q_\lambda,\bm y)}{S^0_k(\bm y)} \;\le\; (1-\lambda)\,\frac{S_k(Q,\bm y)}{S^0_k(\bm y)}+\lambda .
\end{equation}
\end{theorem}
\begin{proof}
By Lemma~\ref{lem:convex}, $S_k(Q_\lambda,\bm y)\le(1-\lambda)S_k(Q,\bm y)+\lambda S_k(F^0,\bm y)$; divide by $S^0_k(\bm y)=S_k(F^0,\bm y)>0$.
\end{proof}
The pinball component is \emph{not} convex in $F$ in general (a quantile of a mixture is not a mixture of quantiles), so the bound is stated
only for CRPS and VS; empirically the tail component behaves similarly. The operator actually shipped is therefore two-parameter:
\begin{equation}\label{eq:action}
  Q_{\gamma,\lambda} \;=\; (1-\lambda)\,Q_\gamma + \lambda\,F^0 ,
\end{equation}
realised with draws: a fraction $1-\lambda$ of systematic-resampled belief draws under weights $\omega^\gamma$, and a fraction $\lambda$ resampled
from M0's own $500$ draws, then shuffled so that draw indices carry no structure.

\subsection{The honesty floor and stepping rule}

Unrestricted, $Q_\gamma$ collapses onto M0's median as $\gamma\to\infty$. We impose:
\begin{enumerate}[leftmargin=*]
\item \textbf{Spread floor.} Reject a candidate if any cell's standard deviation falls below $0.6$ of the honest belief's.
\item \textbf{Effective sample size.} Reject if $\mathrm{ESS}=(\sum_i p_i^2)^{-1}<0.1\,M$ for a pool of $M$ belief draws.
\item \textbf{Stepping.} If the policy's action is rejected, decrease $\gamma$ by $0.25$ until accepted; $\gamma=0$ is always accepted.
\end{enumerate}
For a Gaussian belief equal to M0, the realised single-cell spread ratio is $0.85, 0.74, 0.65, 0.59$ at $\gamma=0.5,1,1.5,2$; with $D=4$ cells it is
$0.94,0.89,0.85,0.83$. So $\gamma=1.5$ is about the largest admissible value on single-cell cards, and after mixing with $\lambda\approx0.3$ the
submitted spread returns to about $0.75$ of M0's. This is a sharp core with the baseline's own tails, not a point forecast. The submission is
concentrated at the random-walk centre, not at any remembered outcome.

\subsection{Why per-card self-optimisation is not used}

An alternative is to choose $\gamma$ per card by minimising a Monte-Carlo estimate of \eqref{eq:risk} under outcomes drawn from the belief. This
is self-confirming: outcomes drawn from $P$ reward sharpening around $P$, so it always selects the largest admissible $\gamma$. It was observed to do
exactly that and was replaced by a fixed, validated action chosen by the policy (Section~\ref{sec:policy}). The evaluation outcomes are in any case
drawn from a robust mixture $(1-r)P+rF^0$ with $r=0.3$, disjoint from the resampling pool.

\section{The numerical prior $P_0$}\label{sec:numerical}

\subsection{Point-in-time state (L1--L3)}

\paragraph{Availability operator.} For a labelled stream $D$, $\Phi(D;\tau)=\{x\in D: t^{\mathrm{avail}}(x)\le\tau\}$. The organizers already
truncate panels and corpora at staging; our loader re-applies $\Phi$ as defence in depth, with the \emph{information cutoff}
$\tau^\star=\min(\tau_{\text{card}},\tau_{\text{cli}})$. Undated documents are excluded because their availability cannot be established.
Non-finite rows are dropped at load.

\paragraph{Steps.} For asset $a$ with gap-cleaned history, $x_{a,t}=v_{a,t}-v_{a,t-1}$ (level) or $x_{a,t}=\log(1+r_{a,t})$ (log return), using
the gap rule of Section~\ref{sec:baseline} on the full history.

\paragraph{State features.} The following are computed at $\tau^\star$:
\begin{align}
  \hat\sigma^2_{E,t} &= \lambda\hat\sigma^2_{E,t-1}+(1-\lambda)x_t^2,\quad \lambda=2^{-1/H},\ H=30 &&\text{(EWMA, uncentred)}\\
  \hat\sigma^2_{300} &= \widehat{\Var}(x_{t-298:t}) &&\text{(M0's window)}\\
  \hat\sigma^2_{L} &= \widehat{\Var}(x_{t-L+1:t}),\quad L=1260 &&\text{(long run)}\\
  \rho_t &= \hat\sigma_{E,t}/\hat\sigma_{L} &&\text{(vol ratio)}\\
  m_t &= \textstyle\sum_{i<63}x_{t-i}\big/(\hat\sigma_L\sqrt{63}) &&\text{(trend $z$-score)}.
\end{align}
The EWMA is initialised by blending in the mean of the first $20$ squared steps with a linearly decaying weight, so the first values are not a
single squared step. The uncentred form is deliberate: the drift of a daily step is negligible relative to its scale, and centring would inject
estimation noise.

\paragraph{Regime (L4).} A deliberately small, soft three-state model on one well-measured quantity, the mean log vol ratio
$\ell=\overline{\log\rho}$ across the card's assets:
\begin{equation}
  \pi_{\text{stress}}=\sigma\bigl(k(\ell-\log 1.35)\bigr),\quad \pi_{\text{calm}}=\sigma\bigl(-k(\ell-\log 0.8)\bigr),\quad
  \pi_{\text{normal}}=\bigl(1-\pi_{\text{stress}}-\pi_{\text{calm}}\bigr)_+ ,
\end{equation}
renormalised, with $k=8$. With about one hundred cards and thin per-regime samples, a high-dimensional latent regime model cannot be estimated
without shrinkage that would defeat it. Soft membership avoids threshold discontinuities. On any failure the belief reduces to one global regime.

\subsection{The variance term structure}

Let $\sigma^2_{\text{now}}=\hat\sigma^2_{E,\tau}$ and let the long-run level be the blend
\begin{equation}
  \sigma^2_{\mathrm{LR}} = (1-\omega)\,\hat\sigma^2_{300}+\omega\,\hat\sigma^2_{L}.
\end{equation}

\begin{proposition}[Aggregated variance under geometric mean reversion]\label{prop:agg}
Suppose one-step conditional variances satisfy $\E_\tau[\sigma^2_{\tau+i}]=\sigma^2_{\mathrm{LR}}+\phi^{\,i-1}(\sigma^2_{\text{now}}-\sigma^2_{\mathrm{LR}})$,
$i\ge1$, and steps are conditionally uncorrelated. Then
\begin{equation}\label{eq:agg}
  V(s) := \Var_\tau\Bigl(\sum_{i=1}^s x_{\tau+i}\Bigr) = s\,\sigma^2_{\mathrm{LR}} + (\sigma^2_{\text{now}}-\sigma^2_{\mathrm{LR}})\,\frac{1-\phi^s}{1-\phi}.
\end{equation}
\end{proposition}
\begin{proof}
With conditionally uncorrelated steps, $\Var_\tau(\sum x_{\tau+i})=\sum_i\E_\tau[\sigma^2_{\tau+i}]$ by the law of total variance (the conditional means
are zero). Summing the geometric series gives $\sum_{i=1}^s\phi^{i-1}=(1-\phi^s)/(1-\phi)$.
\end{proof}
This is the GARCH(1,1) variance term structure with persistence $\phi=\alpha+\beta$. It interpolates between the current EWMA variance at short
horizons and the long-run level at long ones: $V(s)/s\to\sigma^2_{\text{now}}$ as $s\to 1$ and $V(s)/s\to\sigma^2_{\mathrm{LR}}$ as $s\to\infty$. The
implementation floors $V(s)\ge0.05\,s\,\sigma^2_{\mathrm{LR}}$ and takes a running maximum over the grid's step points, so the cumulative path
variance is non-decreasing. We use $\phi=0.985$ per business day (half-life of about $46$ days).

\paragraph{Serial dependence.} Steps are not exactly uncorrelated at the horizon (trend or mean reversion). The empirical variance ratio from
overlapping $s$-sums of demeaned steps,
\begin{equation}
  \widehat{VR}(s)=\frac{\widehat{\Var}\bigl(\sum_{i=1}^s x_{t+i}\bigr)}{s\,\widehat{\Var}(x_t)},
\end{equation}
is noisy, since about $n/s$ independent blocks enter it. We therefore shrink it toward $1$ with a pseudo-count, in empirical-Bayes form:
\begin{equation}\label{eq:vr}
  VR^\star(s) = \mathrm{clip}\Bigl(1+\frac{n/s}{n/s+n_0}\bigl(\widehat{VR}(s)-1\bigr),\,0.6,\,1.6\Bigr),\qquad n_0=40,
\end{equation}
on at most the last $5{,}000$ steps, and $VR^\star\equiv1$ when $n<3s+20$. The final cumulative variance at step count $s$ is
$c^2\,VR^\star(s)\,V(s)$ with a global scale $c$ and per-asset family scale multipliers (Section~\ref{sec:policy}).

\subsection{The centre}

\begin{equation}
  \text{centre}(s) = \text{anchor} + s\bigl[\kappa\,\hat\mu_{300} + \kappa_{\mathrm{LR}}\,\hat\mu_{L}\bigr].
\end{equation}
M0 uses $\kappa=1$. For level targets we use $\kappa=0$ (driftless). For log-return targets the prior is $0.25\hat\mu_{300}+0.75\hat\mu_{L'}$ with
$L'=\max(L,2500)$, and the fitted vintages set the return drift to $0$ as well. The trailing mean has standard error
$\hat\sigma/\sqrt{299}$ per step, so at $s=126$ its contribution to the centre has standard deviation $126\,\hat\sigma/\sqrt{299}\approx7.3\hat\sigma$,
about $0.65$ of the horizon standard deviation $\sqrt{126}\,\hat\sigma\approx11.2\hat\sigma$. For a martingale this is pure noise, and it
enlarges the mean squared error of the centre by roughly $40\%$ of the horizon variance. Removing it improves CRPS on average. It does \emph{not}, however, improve the
ratio loss on its own (Section~\ref{sec:results}), because deviations from M0's centre are penalised convexly.

\subsection{Monthly macro series}

For the monthly CPI/PCE/payroll/unemployment cards the step is a monthly change $\Delta_t$, and inflation and payroll growth are strongly
persistent. With $\tilde\Delta_t$ an EWMA of changes (half-life $6$ months):
\begin{align}
  \text{drift} &= 0.7\,\tilde\Delta_\tau + 0.3\,\overline{\Delta},\\
  \hat\sigma^2_\varepsilon &= \tfrac12\,\mathrm{MAD}(e)^2\cdot1.4826^2 + \tfrac12\,\widehat{\Var}(e_{-60:}),\qquad e_t=\Delta_t-\tilde\Delta_{t-1},\\
  \hat\sigma_\mu &= \sd\bigl(\tilde\Delta_{t}-\tilde\Delta_{t-1}\bigr)_{-60:}\cdot\sqrt6,\\
  V(s) &= s\,\hat\sigma^2_\varepsilon + s^2\,\hat\sigma^2_\mu .
\end{align}
The $s^2$ term represents a slowly moving trend (drift uncertainty grows linearly in the horizon, so its variance grows quadratically). The
MAD blend keeps a single COVID-type print from dominating the innovation scale. The step count $s$ is the M0 month count
(Section~\ref{sec:baseline}), which includes the panel's publication lag.

\subsection{Transfer cards}

On F2 transfer cards the target's own path is withheld: only an early window plus the as-of level is shipped. M0 then estimates its scale from the
early window, which for a formerly managed currency is almost zero. Let $\sigma_{\mathrm{early}}$ be the standard deviation of log steps inside the
early window (gap-dropped) and $\sigma_{\mathrm{proxy}}$ the median, across G10 currencies, of the standard deviation of daily log returns over the last
$300$ observations. We blend in log-variance space:
\begin{equation}
  \log\sigma^2_{\mathrm{rel}} = (1-w)\log\sigma^2_{\mathrm{early}} + w\log\bigl(1.25\,\sigma_{\mathrm{proxy}}\bigr)^2,\qquad w=0.5,
\end{equation}
with an emerging-market premium of $1.25$, and set $\sigma_{\mathrm{level}}=\sigma_{\mathrm{rel}}\cdot|\text{anchor}|$ and $V(s)=s\,\sigma^2_{\mathrm{level}}$.
The geometric blend is appropriate because vol scales multiplicatively and neither source alone is trustworthy. The text layer can then express a
regime (managed versus stressed) through the scenario mixture.

\subsection{Analogues and foundation models}

A same-series analogue engine (Mahalanobis nearest neighbours on vol ratio and $20$/$63$-day trend $z$-scores, estimated on the unit's own history
only, voting on the width) is implemented but disabled, pending a robust walk-forward gain. A time-series foundation model is \emph{not} used: the
artifact policy requires separate approval for any pretrained neural checkpoint, the image must run without a GPU, and the component was never
load-bearing. \code{numerical/foundation.py} is an explicit no-op.

\section{Joint dependence and path-coherent sampling}\label{sec:sampling}

\subsection{Increment correlation}

Let $\Delta$ be the matrix of date-aligned steps for the card's assets. We combine a short and a longer estimator:
\begin{equation}
  C_0 = \tfrac12\,\widehat{\mathrm{Corr}}\bigl(\Delta_{-299:}\bigr) + \tfrac12\,\widehat{\mathrm{Corr}}\bigl(\textstyle\sum_{5}\Delta_{-1260:}\bigr),\qquad
  C = (1-\eta)C_0+\eta\,\widehat{\mathrm{Corr}}(\Delta_{-299:}),\ \eta=0.15,
\end{equation}
where $\sum_5$ denotes overlapping $5$-step sums. These are robust to asynchronous closes (FX fixings versus Treasury closes), which bias daily
correlations toward zero.

\begin{lemma}[Nearest correlation by eigenvalue clipping]\label{lem:nearcorr}
For symmetric $A$ with eigen-decomposition $A=V\Lambda V^\top$, the matrix $\tilde A=D^{-1/2}V\max(\Lambda,\epsilon I)V^\top D^{-1/2}$, with
$D=\diag(V\max(\Lambda,\epsilon I)V^\top)$, is a positive-definite correlation matrix.
\end{lemma}
\begin{proof}
$V\max(\Lambda,\epsilon I)V^\top$ has eigenvalues at least $\epsilon>0$, so it is positive definite. A congruence by the positive diagonal $D^{-1/2}$
preserves positive definiteness and sets the diagonal to one.
\end{proof}
Every correlation matrix produced anywhere in the pipeline passes through this map as a final guard.

\subsection{The generator}

\begin{definition}[Joint belief]
A \code{JointBelief} consists of centres $\mu\in\R^D$; sorted unique step points $p_1<\dots<p_K$; per-asset cumulative variances
$V_a(p_1)\le\dots\le V_a(p_K)$; an increment correlation $C\in\R^{A\times A}$; a scale-mixture dof $\nu$; and scenarios
$\{(\pi_j,\bm b_j,\bm c_j)\}_{j=1}^J$ (weight, per-cell shift, per-cell scale), $\sum_j\pi_j=1$, $\bm c_j>0$.
\end{definition}

For draw $m$:
\begin{enumerate}[leftmargin=*]
\item \textbf{Scenario.} Choose $j_m$ by stratified allocation (exact counts $\lfloor\pi_jN\rfloor$, remainder by largest fractional part at random).
\item \textbf{Mixing.} $W_m=(\nu-2)/\chi^2_\nu$, independent of everything else; $W_m\equiv1$ if $\nu>200$.
\item \textbf{Increments.} $\bm\xi_{m,k}\sim\mathcal N(0,C)$ i.i.d.\ over $k$, and the path
$X_{m,a}(p_k)=\sum_{l\le k}\sqrt{V_a(p_l)-V_a(p_{l-1})}\;\xi_{m,l,a}$ with $V_a(p_0)=0$.
\item \textbf{Value.} For cell $c=(a,h)$ with step count $s_c=p_{k(c)}$:
$Y_{m,c}=\mu_c+b_{j_m,c}+c_{j_m,c}\sqrt{W_m}\,X_{m,a}(s_c)$, then the soft floor where one applies.
\end{enumerate}
Draws are generated in antithetic pairs $(\bm\xi,-\bm\xi)$ sharing $W$ and the scenario label.

\begin{proposition}[Covariance of the base layer]\label{prop:pathcov}
For the base scenario ($\bm b=0,\bm c=1$) and $\nu>2$, for cells $c=(a,s)$ and $c'=(a',s')$,
\begin{equation}
  \Cov(Y_c,Y_{c'}) \;=\; C_{aa'}\sum_{l:\,p_l\le\min(s,s')}\sqrt{\Delta V_a(p_l)\,\Delta V_{a'}(p_l)},\qquad \Delta V(p_l)=V(p_l)-V(p_{l-1}).
\end{equation}
In particular $\Var(Y_c)=V_a(s)$, and for one asset $\Cov(Y_{(a,s)},Y_{(a,s')})=V_a(\min(s,s'))$: the path structure of
Proposition~\ref{prop:m0path}, with the random-walk variance replaced by the term structure.
\end{proposition}
\begin{proof}
$\E W=1$ (Lemma~\ref{lem:tmix}), and $W$ is independent of $\bm\xi$, so $\Cov(\sqrt W X_{a}(s),\sqrt W X_{a'}(s'))=\E[W]\,\Cov(X_a(s),X_{a'}(s'))$.
Increments at different $l$ are independent with zero mean, so only common indices $l\le\min$ contribute, each
$\sqrt{\Delta V_a\Delta V_{a'}}\,\E[\xi_{l,a}\xi_{l,a'}]=\sqrt{\Delta V_a\Delta V_{a'}}\,C_{aa'}$.
\end{proof}

\begin{corollary}[Validity]
The generator always produces a valid joint law: its covariance is a sum over $l$ of $D_l C_{\mathcal A} D_l$ terms with diagonal $D_l\succeq0$,
each positive semidefinite. No matrix is ever edited element-wise.
\end{corollary}

\begin{lemma}[Variance-preserving $t$ mixing]\label{lem:tmix}
If $\chi^2\sim\chi^2_\nu$ with $\nu>2$ then $\E[(\nu-2)/\chi^2]=1$. Hence $\sqrt W\,\bm Z$ has the covariance of $\bm Z$, is multivariate Student-$t$ with
$\nu$ degrees of freedom, and has tail dependence across cells.
\end{lemma}
\begin{proof}
$\E[1/\chi^2_\nu]=1/(\nu-2)$ for $\nu>2$, since $\int_0^\infty x^{-1}\frac{x^{\nu/2-1}e^{-x/2}}{2^{\nu/2}\Gamma(\nu/2)}dx=\frac{2^{\nu/2-1}\Gamma(\nu/2-1)}{2^{\nu/2}\Gamma(\nu/2)}=\frac1{\nu-2}$.
\end{proof}
Mixing preserves the variance forecast exactly: the choice of tail shape does not bias the width. The walk-forward selected $\nu\approx\infty$
(Gaussian) for the fitted vintages. Fat tails were rejected by the ratio metric (Section~\ref{sec:results}).

\subsection{The soft lower bound for yields}

Gaussian walks from $0.1\%$ put real mass on yields that do not print. For UST cells we apply
\begin{equation}
  f(y)=\begin{cases} y & y\ge b\\ b-k\log\bigl(1+(b-y)/k\bigr) & y<b\end{cases},\qquad b=f_0+k,\ f_0=-0.10,\ k=0.25,
\end{equation}
with a hard stop at $f_0-5k$.

\begin{proposition}
$f$ is continuous, continuously differentiable, strictly increasing, the identity on $[b,\infty)$, and $f'(b^-)=1$.
\end{proposition}
\begin{proof}
For $y<b$, $f'(y)=k\cdot\frac{1/k}{1+(b-y)/k}=\frac{1}{1+(b-y)/k}\in(0,1)$, which tends to $1$ as $y\uparrow b$, and $f(b^-)=b$.
\end{proof}
Because $f$ is monotone it maps quantiles to quantiles and leaves the copula untouched. The upper tail, which matters for hiking episodes, is unaffected.

\subsection{Text-driven dependence (L19)}

When the text supports a stress regime on a multi-asset card, the increment correlation moves convexly toward a one-factor stress structure
$C^{\text{stress}}_{ij}=s_is_j\,r$ for $i\ne j$, where $s_i$ is the sign of asset $i$'s loading on the first principal component and
$r=\min\{0.95,\ \overline{|C_{ij}|}+\tfrac12(1-\overline{|C_{ij}|})\}$:
\begin{equation}
  C_1 = (1-w)\,C + w\,C^{\text{stress}},\qquad w=\min\bigl(0.5,\ 0.5\,g\cdot\text{dependence strength}\bigr).
\end{equation}
\begin{proposition}\label{prop:psd}
If $C,C^{\text{stress}}\succeq0$ then $C_1\succeq0$ for every $w\in[0,1]$, and $C_1=C$ at $g=0$.
\end{proposition}
\begin{proof}
For $v\in\R^A$, $v^\top C_1v=(1-w)v^\top Cv+w\,v^\top C^{\text{stress}}v\ge0$; the positive-semidefinite matrices form a convex cone. The matrix
$C^{\text{stress}}=r\,\bm s\bm s^\top+(1-r)I$ is positive semidefinite for $r\in[0,1]$.
\end{proof}

\section{The text branch: from documents to a structured signal}\label{sec:text}

The text branch runs independently of the numerical branch and communicates with it through one narrow interface, the triple $(Z,g,\kappa)$. It
uses no neural weights: passages, dictionaries and a BM25 index are built at run time from \emph{this unit's} corpus only, so nothing can leak
across units.

\subsection{Events (L8)}

Each document is split into paragraph blocks, packed greedily into passages of at most $900$ characters (longer paragraphs are cut). A passage
$e$ is kept if it contains any market vocabulary. With counts $h_e,d_e$ of hawkish and dovish phrases, $u_e$ of stress terms and $\tau_e$ of
explicit tail terms (all from static dictionaries),
\begin{equation}
  \text{tone}_e=\frac{h_e-d_e}{h_e+d_e+2}\in(-1,1),\qquad \text{stress}_e=\min(1,u_e/6).
\end{equation}
The $+2$ is a Laplace-type pseudo-count that keeps a single phrase from saturating the tone. Assets touched are those whose vocabulary appears.

\subsection{Deduplication (L9)}

Let $\mathcal S_e$ be the set of hashed word $5$-grams of $e$. Two passages from different documents are merged iff
\begin{equation}
  J(e,e')=\frac{|\mathcal S_e\cap\mathcal S_{e'}|}{|\mathcal S_e\cup\mathcal S_{e'}|}\ge\theta=0.55
  \quad\text{and}\quad |t_e-t_{e'}|\le 5\ \text{days}.
\end{equation}
Clusters are the connected components of this relation (union-find over a time-sorted sweep, $O(n\cdot w)$ for window population $w$). A
cluster's corroboration is the number $c$ of distinct documents in it, and its novelty weight is $\nu=1/(1+\log c)$. Requiring both semantic and
temporal proximity keeps a topic recurring months apart as two events. Without deduplication, $k$ wire copies of one statement would look like
$k$ independent confirmations and inflate the confidence the gate relies on.

\subsection{State-conditioned query and retrieval (L10--L11)}

The query is built from the numerical state, not from "recent news": the vocabulary of the target assets, stress and tail vocabulary when
$\pi_{\text{stress}}>0.4$, and content words from the card's own authored title, description and notes (a legitimate input file). BM25 with
$k_1=1.4$, $b=0.75$ scores unigram terms:
\begin{equation}
  \mathrm{BM25}(e)=\sum_{t\in q}\mathrm{idf}(t)\,\frac{f_{t,e}(k_1+1)}{f_{t,e}+k_1\bigl(1-b+b\,|e|/\overline{|e|}\bigr)},\quad
  \mathrm{idf}(t)=\log\Bigl(1+\frac{N-n_t+0.5}{n_t+0.5}\Bigr),
\end{equation}
plus $1.5$ per matched multi-word phrase. With $\mathrm{rel}_e=\mathrm{BM25}(e)/\max_{e'}\mathrm{BM25}(e')$, the final score merges the five
retrieval channels of the reference designs into one:
\begin{equation}\label{eq:retr}
  \mathrm{score}(e)=\mathrm{rel}_e+\sqrt{\mathrm{rel}_e}\Bigl[0.6\,e^{-\lambda\,\mathrm{age}_e}+0.8\,\mathrm{match}_e+0.3\,\nu_e+0.4\,w_{\mathrm{type}(e)}\Bigr],
\end{equation}
with half-life $45$ days, $\mathrm{match}_e$ the fraction of target assets touched, and source-type weights (FOMC statement $1.0$, testimony and
minutes $0.8$, speeches $0.6$, \ldots). The factor $\sqrt{\mathrm{rel}_e}$ was introduced after a test exposed a real flaw: with purely additive
terms, a recent passage of zero relevance outranked a relevant older one. Now recency, asset match, novelty and source type can re-order relevant
passages but cannot lift an irrelevant one, since the score is $0$ when BM25 is $0$.

\subsection{Diversity and budget (L12)}

Maximal marginal relevance selects up to $12$ passages within $14{,}000$ characters:
\begin{equation}
  e^\star=\argmax_{e\notin\mathcal C}\ \Bigl[0.7\,\mathrm{score}(e)-0.3\cdot3\max_{e'\in\mathcal C}J_1(e,e')\Bigr],
\end{equation}
with unigram Jaccard $J_1$. Each passage is compressed to whole sentences within its share of the budget. Redundancy is a property of a
\emph{set}, so this is where the redundancy term of the reference retrieval score belongs.

\subsection{One structured call (L13)}

The House model receives the as-of date, the card's authored notes, the quote convention of every target (for example "EUR: US dollars per
1 EUR; up = EUR strengthens", "JPY: JPY per 1 US dollar; up = JPY weakens"), the horizons, a compact numeric state (last value, trend $z$, vol ratio,
regime probabilities) and the compressed evidence with document identifiers. It is instructed to use only that information set and to ignore any
recognition of the episode, and it is asked for
\begin{multline}
  Z=\Bigl(\{p^{\uparrow}_a,\ m_a,\ v_a,\ \text{tail}_a\}_{a\in\mathcal A},\ \{(\text{name}_s,\pi_s,\bm d_s,m_s,v_s)\}_{s\le3},\\
  \text{persistence},\ \text{regime},\ C,\ \nu,\ \text{established},\ \text{cited}\Bigr),
\end{multline}
with $p^\uparrow_a\in[0,1]$ the probability of ending above the numeric centre, $m_a\in\{$none, weak, moderate, strong, extreme$\}$,
$v_a\in\{$lower, same, higher, much\_higher$\}$ and $\text{tail}_a\in\{$none, up, down, both$\}$. The model is never asked for a number for
the market.

\paragraph{Call discipline.} Temperature $0$, a pinned seed, $n=1$, thinking disabled through \code{chat\_template\_kwargs} (thinking tokens
would consume the $4{,}000$-token cap), \code{max\_tokens}$\,=1800$, and no tools. Our own cap is $3$ calls per card against the allowance of $25$,
with at most one retry after a parse failure. An HTTP $400$/$422$ triggers one resend with the bare OpenAI body, for a deployment that rejects the
optional fields. Any other failure, and any budget or deadline exhaustion, yields \emph{no} signal.

\subsection{Schema coercion as a one-sided operation}

\begin{definition}
A coercion $\mathcal K$ from raw model output to $Z$ is \emph{one-sided} if every malformed field maps to its neutral value: $p^\uparrow=\tfrac12$,
$m=$ none, $v=$ same, tail $=$ none, confidence $0$, established $=$ false.
\end{definition}
The implemented coercion is one-sided. Numbers are accepted as probabilities in $[0,1]$ or percentages in $(1,100]$, and anything else maps to
the default. Only a JSON boolean \code{true} (or the literal string \code{"true"}) counts as established. Two violations of one-sidedness were found
and fixed by adversarial tests: the string \code{"yes"} had been read as established, which is the one flag that permits narrowing, and a
confidence of $250$ had been clipped to the \emph{maximum} $1$ instead of being treated as unusable. JSON is extracted with a
string-and-escape-aware brace matcher after stripping \code{<think>} blocks and code fences, tolerating trailing commas.

\subsection{Lexicon fallback}

Without a usable model answer, a deliberately weak signal is derived from the retrieved passages: widening only when the mean stress exceeds
$0.35$ ("higher") or $0.6$ ("much higher"), a direction only for rates and only from Federal Reserve tone ($p^\uparrow=0.5+0.08\,\overline{\text{tone}}$),
and confidence at most $0.4$. The gate additionally caps a lexicon-sourced signal (Section~\ref{sec:fusion}).

\section{Fusion: gating, contradiction and the distribution update}\label{sec:fusion}

\subsection{Contradiction (L16)}

For asset $a$ with text view $t_a=2p^\uparrow_a-1\in[-1,1]$, numeric trend $z$-score $m_a$ and model confidence $C$,
\begin{equation}
  \kappa_a = C\,|t_a|\;\ind\{\sgn t_a\ne\sgn m_a\}\;\min\bigl(1,|m_a|/1.5\bigr),\qquad \kappa=\overline{\kappa_a}\in[0,1].
\end{equation}
A confident text view against a clear numeric trend counts most, and an unsure view or a flat trend contradicts little. Contradiction is never
resolved in favour of one source: it lowers the gate and widens the distribution.

\subsection{The information gate (L15)}

\begin{equation}\label{eq:gate}
  g=\min\Bigl\{g_{\max},\ \sigma\bigl(w_0+\delta_f+|w_1|\,q+|w_2|\,C+|w_3|\log(1+c_{\max})-|w_4|\,\kappa\bigr)\cdot\theta_f\cdot a\Bigr\},
\end{equation}
where $q\in[0,1]$ is the retrieval quality (the top-$3$ scores relative to the maximum), $c_{\max}$ the largest corroboration among retrieved
passages, $\delta_f$ the family intercept shift, $\theta_f$ the family text strength, $a$ the policy's action multiplier and $g_{\max}=0.85$. The
defaults are $w_0=-1.2$, $w_1=1.1$, $w_2=2.2$, $w_3=0.35$, $w_4=1.4$, with $(\delta,\theta)=(-0.3,0.8),(0.4,1.0),(0,0.8),(0.2,0.9)$ for F1--F4.
Hard rules: a neutral signal or an empty retrieval set gives $g=0$ \emph{exactly}; a lexicon-sourced signal is capped at $g\le0.25$.

\begin{proposition}[Monotonicity]\label{prop:gate}
$g$ is non-decreasing in $q$, $C$ and $c_{\max}$, non-increasing in $\kappa$, and $g\in[0,g_{\max}]$.
\end{proposition}
\begin{proof}
The logistic function is increasing; the argument is non-decreasing in $q,C,c_{\max}$ and non-increasing in $\kappa$, because the coefficients
enter through absolute values. Multiplication by $\theta_fa\ge0$ and the cap $\min\{g_{\max},\cdot\}$ preserve monotonicity, and all factors are
non-negative.
\end{proof}
The absolute values replace the softplus reparametrisation of the design document: an unconstrained fit on a handful of cards could otherwise learn
a sign-flipped relation that happens to fit noise. The weights are priors, not fitted. With no House endpoint off-platform there is no dataset on
which $w$ could be estimated without look-ahead (Section~\ref{sec:limitations}).

\subsection{The distribution update (L18)}

Let $\sigma_c$ be the base standard deviation of cell $c$ and $m(\cdot)$ the \emph{semantic-to-empirical} magnitude map, anchored on quantiles of
$|z|$ for standardised moves: $m(\text{weak})=0.674$, $m(\text{moderate})=1.150$, $m(\text{strong})=1.645$, $m(\text{extreme})=2.326$, the
$50$/$75$/$90$/$98$\% points of $|\mathcal N(0,1)|$. When the in-unit residual bank is enabled these are averaged half-and-half with the unit's own
empirical $|z|$ quantiles. Persistence enters as
\begin{equation}
  f_{\mathrm{pers}}(s)=\begin{cases}\sqrt{\min(1,21/s)} & \text{transient}\\ \sqrt{\min(1,63/s)} & \text{medium}\\ 1 & \text{persistent.}\end{cases}
\end{equation}

\paragraph{Mechanism A (tilt).} With $\alpha=0.25\,\theta_f$, $\beta=0.30$, $\lambda_\kappa=0.15$ and the volatility codes
$v(\text{lower})=-\tfrac12$, $v(\text{same})=0$, $v(\text{higher})=1$, $v(\text{much\_higher})=2$,
\begin{align}
  \mu_{1,c} &= \mu_{0,c}+g\,\alpha\,(2p^\uparrow_{a}-1)\,m(m_a)\,f_{\mathrm{pers}}(s_c)\,\sigma_c, \label{eq:tilt-mu}\\
  \sigma_{1,c} &= \sigma_{0,c}\,\exp\Bigl(\min\bigl\{g\beta\,v^\star_a+\lambda_\kappa\kappa\sqrt g,\ \log 2\bigr\}\Bigr), \label{eq:tilt-s}
\end{align}
where $v^\star_a=\max(v_a,0)$ unless the driver is \emph{established}. Inferred tone may move the centre but may not narrow.

\begin{proposition}[Validity and exact reduction]\label{prop:tilt}
(i) $\sigma_{1,c}>0$ whenever $\sigma_{0,c}>0$. (ii) If $g=0$ then $(\mu_1,\sigma_1)=(\mu_0,\sigma_0)$ exactly. (iii) $\sigma_{1,c}\le2\,\sigma_{0,c}$.
(iv) $|\mu_{1,c}-\mu_{0,c}|\le g\,\alpha\,m(\text{extreme})\,\sigma_c\le 0.85\cdot0.25\cdot2.326\,\sigma_c\approx0.49\,\sigma_c$.
\end{proposition}
\begin{proof}
(i) $\exp(\cdot)>0$. (ii) Every term in the exponent and the shift carries a factor $g$ or $\sqrt g$. (iii) The exponent is capped at $\log2$.
(iv) $|2p-1|\le1$, $f_{\mathrm{pers}}\le1$, $\theta_f\le1$ and $g\le g_{\max}$.
\end{proof}
The multiplicative form excludes an entire class of invalid states: an additive $\sigma_1=\sigma_0+\Delta$ can become non-positive under a confident
but wrong signal. The cap and the reduced $\alpha$ (from $0.35$) are the result of the audit in Section~\ref{sec:engineering}. Under the convex ratio loss a
centre shift of about $0.57\sigma$, the old worst case, is too costly when the text is wrong.

\paragraph{Mechanism B (mixture).} When the model describes at least two branching scenarios, or $\kappa>\kappa^\star=0.55$, the belief becomes
\begin{equation}\label{eq:mixture}
  P_1=(1-g)\,P_0^{(\sigma_1)}+g\sum_{s}\pi_s\,P_s,
\end{equation}
where component $P_s$ shifts cell $c$ by $d_{s,a}\,m(m_s)\,f_{\mathrm{pers}}(s_c)\,\sigma_c\min(1,1.5\alpha)$ and scales it by
$\min\{e^{\beta v^\star_s},2\}\,\sigma_{1,c}$,
where $P_0^{(\sigma_1)}$ is the prior with the widened scale.

\begin{proposition}\label{prop:mix}
For $g\in[0,1]$ and densities $P_0,P_s$, \eqref{eq:mixture} is a probability density, and $P_1=P_0^{(\sigma_1)}=P_0$ at $g=0$.
\end{proposition}
\begin{proof}
It is a non-negative combination with weights $(1-g)+g\sum_s\pi_s=1$, because the coerced scenario probabilities are renormalised to sum to one.
Integrating gives $1$ by linearity. At $g=0$ only the first term remains, and $\sigma_1=\sigma_0$ by Proposition~\ref{prop:tilt}(ii).
\end{proof}
Each draw selects one scenario applied to \emph{all} cells, so a scenario is a coherent joint story ("the central bank blinks") rather than a
per-cell coin flip. This is what makes text useful for the variogram term.

\paragraph{Event-triggered one-sided tail (L20, F4-weighted).} If any asset's tail flag is set, a component of weight
\begin{equation}
  w_{\mathrm{tail}}=\min\bigl\{0.25,\ g\cdot0.12\cdot\theta^{\mathrm{tail}}_f\cdot(1\text{ if strong or established else }\tfrac12)\bigr\}
\end{equation}
is added, with scale $2.2\,\sigma_1$ and shift $\pm m(\text{strong})\,\sigma_c$ toward the flagged side (no shift for "both"). The other weights are
multiplied by $1-w_{\mathrm{tail}}$. This targets the $1\%$ and $99\%$ pinball points directly instead of widening the whole distribution, which
would cost CRPS on ordinary days.

\subsection{The semantic-to-numeric map as a bottleneck}

The LLM's output enters the distribution only through the finite set of categorical fields above and the bounded probabilities $p^\uparrow$, $\pi_s$
and $C$. Numbers are produced by the panel ($\sigma_c$) and by fixed maps ($m$, $f_{\mathrm{pers}}$, $v$). This low-dimensional bottleneck is what
makes the text component auditable and bounded (Proposition~\ref{prop:tilt}), and it is the defence against an open model's miscalibrated
numerical intuition.

\section{Tails, shrinkage and calibration}\label{sec:tails}

\subsection{The in-unit residual bank}

For past cutoffs $t$ of the unit's own history (stride $10$ over the last $12$ years, and only $t+s\le$ last observation $\le\tau$), the same
variance formula gives $V(s\mid t)$ from $\sigma^2_{E,t}$, $\hat\sigma^2_{300,t}$ and $\hat\sigma^2_{L,t}$ as vectorised rolling series, and the
standardised error is
\begin{equation}
  z_t(s)=\frac{\sum_{i=1}^s x_{t+i}}{\sqrt{c^2\,VR^\star(s)\,V(s\mid t)}} .
\end{equation}
Every quantity is computable at $\tau$: this is a walk-forward \emph{inside} the container, needing no cross-unit or post-cutoff data.

\subsection{Empirical-Bayes shrinkage with effective sample size}

Overlapping windows at stride $k$ carry fewer than $n$ independent observations.

\begin{lemma}[Effective size under overlap]
If $s$-step sums are sampled at stride $k<s$ from a series of independent steps, adjacent sums share a fraction $1-k/s$ of their steps. The
number of non-overlapping blocks, which bounds the independent information, is about $n\,k/s$.
\end{lemma}
\begin{proof}
Over a span of $nk$ steps, non-overlapping $s$-blocks number $nk/s$. Any sum is a linear combination of these blocks and of partial blocks, and
the variance estimate cannot use more independent components than there are blocks.
\end{proof}
We use $n_{\mathrm{eff}}=n\min(1,k/s)$. The scale bias and its shrinkage toward the global level $1$ are
\begin{equation}
  \hat b=\sqrt{\mathrm{RMS}(z)\cdot1.4826\,\mathrm{MAD}(z)},\qquad
  \log b^\star=\frac{n_{\mathrm{eff}}}{n_{\mathrm{eff}}+n_0}\log\hat b,\quad n_0=60,\qquad b^\star\in[0.75,1.35],
\end{equation}
the geometric mean of two scale estimators that are both consistent under normality. For tail thickness, with sample excess kurtosis $\hat k$ of
$z$, the local Student dof is $\hat\nu=4+6/\hat k$, because a $t_\nu$ has excess kurtosis $6/(\nu-4)$. It is shrunk in $1/\nu$, the natural
parameter for tail weight:
\begin{equation}
  \frac1{\nu^\star}=\frac{n_{\mathrm{eff}}}{n_{\mathrm{eff}}+n_0'}\cdot\frac1{\hat\nu}+\frac{n_0'}{n_{\mathrm{eff}}+n_0'}\cdot\frac1{\nu_{\mathrm{prior}}},\qquad n_0'=120 .
\end{equation}

\begin{remark}[Why shrinkage has this form]
If $\log\hat b\sim\mathcal N(\log b,\tau^2/n_{\mathrm{eff}})$ and the prior is $\log b\sim\mathcal N(0,\tau^2/n_0)$, the posterior mean is
exactly $\frac{n_{\mathrm{eff}}}{n_{\mathrm{eff}}+n_0}\log\hat b$: the James--Stein/empirical-Bayes argument, applied to a scale parameter.
\end{remark}

\paragraph{Status.} The bank is implemented and tested but \emph{disabled}. In walk-forward tuning it did not deliver a robust out-of-sample gain
(train $0.948$ against $0.945$ for the base; test $0.964$ against $0.967$), which is within noise. The pre-registered rule is that a component
ships only on a robust gain.

\subsection{Probabilistic calibration}

If realised PITs $u=F(y)$ have cdf $G$, then $F'=G\circ F$ is calibrated. For draws, a draw at forecast level $u$ is moved to the forecast's
own quantile at level $G^{-1}(u)$, per cell, using the draws' ranks with one common $u$ per draw.

\begin{proposition}[Calibration leaves the copula unchanged]
Applying a strictly increasing map $T_c$ to each cell separately leaves the copula of the joint draws unchanged.
\end{proposition}
\begin{proof}
The copula is the joint law of the ranks, and strictly increasing maps preserve ranks (Sklar's theorem).
\end{proof}
$G^{-1}$ is estimated on $17$ knots by pool-adjacent-violators isotonic regression of the empirical PIT quantiles, then shrunk toward the identity
with weight $n/(n+200)$, per asset class $\times$ horizon bucket.

\paragraph{Status: rejected.} Maps fitted on outcomes before 2015 \emph{worsened} the 2015--2024 walk-forward from $0.964$ to $0.977$, helping on
only $48\%$ of cards. The reason is instructive. The PIT histograms show mild miscalibration of the \emph{honest} belief, but the submission is
not the honest belief. The score-aware operator already reshapes it for the ratio loss, and recalibrating the input to that operator does not carry
through. Calibration is disabled, and the maps are kept in \code{calibration\_vintages.json} for transparency.

\section{The adaptive policy: reinforcement learning under noisy rewards}\label{sec:policy}

\subsection{What the policy is allowed to decide}

The statistical backbone is fixed. The policy only decides how far the submission moves from the honest belief toward the metric-optimal act,
through the discrete action set
\begin{equation}
  \mathcal A=\{(\gamma,\lambda):\gamma\in\{0,0.5,0.8,1,1.25,1.5,1.75,2\},\ \lambda\in\{0,0.1,0.2,0.3,0.4\}\},\qquad |\mathcal A|=40,
\end{equation}
applied through \eqref{eq:action}. Text strength is \emph{not} an action. The historical walk-forward has no corpus, so there is no data from which
to learn it, and a text action would be learned from nothing. It stays with the family policy and the gate.

\subsection{Why this is full-information contextual learning, not PPO}

A card is a one-step episode: the action does not change future states, only the loss on this card. That is a contextual decision problem. In
the offline walk-forward the outcome $\bm y$ is observed after the episode, and the loss of \emph{every} action on the same belief pool is
computable. The feedback is therefore full-information, not bandit: there is no exploration problem. The problem reduces to cost-sensitive
contextual learning:
\begin{equation}
  \pi^\star(x)=\argmin_{a\in\mathcal A}\ \E\bigl[\Snorm(Q_a,\bm Y)\,\big|\,x\bigr].
\end{equation}

\subsection{Variance-controlled reward}

Financial outcomes are dominated by noise: a four-sigma surprise scores badly under every action. Regressing the raw loss on the context would
mostly learn card difficulty. We instead use the \emph{advantage} over a default action $a_0$ on the same card and the same outcome,
\begin{equation}
  r_a = \Snorm(Q_{a_0},\bm y)-\Snorm(Q_a,\bm y),
\end{equation}
which cancels the common card-level component. Formally, if $\Snorm(Q_a,\bm y)=\beta(\bm y)+\delta_a(x,\bm y)$ with a large common term $\beta$,
then $\Var(r_a)=\Var(\delta_{a_0}-\delta_a)\ll\Var(\beta)$. This is the counterfactual incremental reward of the design notes, realised exactly
because every action is scored.

\subsection{Context}

Sixteen features, all available at inference time:
\begin{equation}
  x=\bigl(1,\ \overline{\log\rho},\ \overline{\min(|m|,4)},\ \log(\bar s/21),\ \ind_{\text{single}},\ \ind_{\text{log\_return}},\
  \{\text{class shares}\}_{5},\ \{\ind_{\text{family}}\}_{4},\ g_{\text{text}}\bigr),
\end{equation}
where $g_{\text{text}}$ is $0$ at decision time (the policy is chosen before the text runs).

\subsection{Per-action Bayesian ridge and pessimistic selection}

For each non-default action $a$ we fit $r_a\approx x^\top\theta_a$ with a Gaussian prior $\theta_a\sim\mathcal N(0,\sigma^2\Lambda^{-1})$,
$\Lambda=\lambda_rI$:
\begin{equation}
  A_a=\lambda_rI+X^\top X,\qquad \hat\theta_a=A_a^{-1}X^\top r_a,\qquad \Var(x^\top\theta_a\mid\mathcal D)=\hat\sigma^2_a\,x^\top A_a^{-1}x,
\end{equation}
with $\hat\sigma^2_a$ the residual variance and $\lambda_r=100$. Selection is \emph{pessimistic}:
\begin{equation}\label{eq:lcb}
  \hat a(x)=\begin{cases}\argmax_a\mathrm{LCB}_a(x) & \max_a\mathrm{LCB}_a(x)>0\\ a_0 & \text{otherwise,}\end{cases}\qquad
  \mathrm{LCB}_a(x)=x^\top\hat\theta_a-k\,\hat\sigma_a\sqrt{x^\top A_a^{-1}x}.
\end{equation}

\begin{proposition}[Safe improvement]\label{prop:safe}
If the posterior is correctly specified, then conditional on $x$ and for each $a$, $\PP\bigl(\E[r_a\mid x]<0\ \text{and}\ \hat a(x)=a\bigr)\le\Phi(-k)$.
Deviations from the default therefore rarely make the expected loss worse; with $k=2$, the per-action rate is at most $2.3\%$.
\end{proposition}
\begin{proof}
$\hat a(x)=a$ requires $\mathrm{LCB}_a(x)>0$, that is $x^\top\hat\theta_a>k\,\mathrm{sd}$. Under the posterior, $\E[r_a\mid x]=x^\top\theta_a$ is
Gaussian with mean $x^\top\hat\theta_a$ and that standard deviation, so the posterior probability that it is negative is at most $\Phi(-k)$.
\end{proof}
This is the spirit of safe policy improvement with baseline bootstrapping: deviate from a trusted default only where the data support it with
margin. \code{LinearArm} also provides Thompson sampling and LinUCB-style online updates, but they are not used in production, because offline
feedback is full-information.

\subsection{Evidence that pessimism is necessary}

Trained on cards whose outcome is known before 2015 and evaluated on cards from 2015 onwards:
\begin{center}
\begin{tabular}{lcccc}
\toprule
ridge $\lambda_r$ & pessimism $k$ & test, default & test, policy & share deviating\\
\midrule
20 & 0 & 0.9665 & 0.9689 & 96\%\\
20 & 2 & 0.9665 & \textbf{0.9636} & 36\%\\
100 & 0 & 0.9665 & 0.9704 & 94\%\\
100 & 1 & 0.9665 & 0.9664 & 70\%\\
100 & 2 & 0.9665 & \textbf{0.9631} & 31\%\\
400 & 2 & 0.9665 & 0.9649 & 17\%\\
\bottomrule
\end{tabular}
\end{center}
With no pessimism the policy improves in-sample (train about $0.936$ against $0.945$) and \emph{worsens} out of sample. That is overfitting to noise, as
one should expect with low signal-to-noise. Only the pessimistic variants beat the default out of sample, and modestly: $0.9665\to0.963$. The
shipped policy uses $\lambda_r=100$, $k=2$, fitted per vintage (Section~\ref{sec:validation}). For F2 and F4 in the latest vintage the policy is
switched off, because those families use a different default action. A policy trained against the generic default would steer toward the wrong
reference.

\subsection{Family policy}

A family changes only a handful of scalars (gate intercept shift, text strength, tail strength, dependence strength, a width multiplier, draw
count), never the backbone. F4 uses at least $4{,}000$ draws, because the $1\%$ quantile from $200$ draws rests on two order statistics.

\section{Validation methodology}\label{sec:validation}

\subsection{Why the practice leaderboard cannot be the primary signal}

Practice cards share the same underlying series and are truncated only at their own as-of dates. The organizers measured that for $75$ of $103$
practice units every realised value appears, to full precision, in a sibling unit's published panel. Any improvement on the practice board is
therefore confounded with recovering leaked values, which we never do. The primary signal is a self-constructed walk-forward; the practice board
is a check of admissibility and, once per upload, a check of fidelity (Section~\ref{sec:results}).

\subsection{History reconstruction}

The full panels are not shipped. We rebuild them as the union of the practice units' truncated copies (tuning on published practice data is
permitted). The consistency audit found $0$ conflicting $(\text{asset},\text{date})$ values once the worked exemplar was excluded: that unit ships an
illustrative non-FRED 2024 panel, which initially produced $496$ conflicts of up to $75$\,bp. Coverage: Treasuries 2000-01-03 to 2024-12-18, G10 FX to
2024-10-31, factors to 2024-05-31. No post-2024 data were downloaded: the sealed window is H2~2025 to Q2~2026, and tuning on data after an evaluation
task's cutoff would breach the artifact policy.

\subsection{Synthetic cards}

Thirteen templates mirror the practice roster's shapes: single-cell rates and FX at long ($126/189$), mid ($63/64$) and short ($21$) horizons;
factor log returns; UST curves ($4$ tenors $\times$ $\{63,126\}$ or $\{21,63\}$); G10 baskets; mixed rates-and-FX; factor pairs. As-of dates are every $21$
business days from 2003, with realised values read from the $h$-th panel row after the as-of (or the sum of $\log(1+r)$ over it). This gives
$3{,}896$ cards, each recording the last date its outcome needs (\code{target\_end}).

\subsection{Production-exact evaluation}

The lab calls the production function \code{forecast\_from\_pit} with an in-memory point-in-time slice, so the formulas are never forked between
research and production. Each card is scored with the exact M0 replica and the exact metric, which is the leaderboard statistic itself. For every
card the runner also scores every action $(\gamma,\lambda)$ on the same belief pool (full information), rejecting actions that break the spread
floor.

\subsection{Vintaged artifacts: out-of-sample by construction}

\begin{definition}[Vintage]
A vintage $v$ is a set of fitted overrides with a \code{fit\_cutoff} $t_v$ equal to the last date of any outcome used to fit or select it. A card
with as-of $\tau$ uses the vintage $v^\star(\tau)=\argmax\{t_v: t_v<\tau\}$, or built-in priors if none exists.
\end{definition}

\begin{proposition}[No look-ahead through parameters]
Under this rule, every parameter used to forecast a card was a measurable function of outcomes dated strictly before the card's as-of.
\end{proposition}
\begin{proof}
By construction the fitting set of $v$ contains only cards with \code{target\_end}$\,\le t_v$, and $t_v<\tau$ for the selected vintage.
\end{proof}
The shipped vintages are $\{$2008-12-31, 2014-12-31, 2024-12-31$\}$. Cards dated 2003--2008 use theory priors ($\gamma=0.5$ from the simulation of
Section~\ref{sec:ratio}, $\lambda=0.4$), cards dated 2009--2014 the 2008 vintage, cards dated 2015--2024 the 2014 vintage, and the sealed 2025+ cards
the 2024 vintage. Consequently the whole 2003--2024 walk-forward, run with the production vintage rule, is an out-of-sample evaluation of the
system that ships. The one exception is the variant-level choice of the belief model, which was made once on the pre-2015 split and is within noise.

\subsection{Uncertainty: a block bootstrap over quarters}

Cards with nearby as-of dates share outcome windows, so resampling cards independently would understate uncertainty. We resample \emph{calendar
quarters} of as-of dates with replacement and recompute the mean. For $Q$ quarters with totals $(T_q,n_q)$ and resampled indices $q_1,\dots,q_Q$,
\begin{equation}
  \bar S^{*}=\frac{\sum_{i}T_{q_i}}{\sum_in_{q_i}},\qquad \text{CI}=\bigl[\mathrm{quantile}_{0.025}(\bar S^*),\ \mathrm{quantile}_{0.975}(\bar S^*)\bigr],
\end{equation}
from $400$ replications. This is the moving-block idea with blocks aligned to the calendar, which respects the dependence induced by overlapping
horizons.

\subsection{The ablation ladder and multiple testing}

The rungs and their order were fixed before results were inspected: M0 (replica), M1 (driftless random walk at M0's width), M2 (variance term
structure), M3 (plus analogues), M4--M8 (text rungs, on practice units: text, gate, contradiction, distribution update, dependence), M9 (fat tails
and residual bank), M10 (calibration), M11 (full), M12 (stress slices). Tuning was one-factor-at-a-time around a base over a fixed list of $14$
variants, selected on train and \emph{reported} on test. No open-ended search was run.

\section{Results}\label{sec:results}

All walk-forward numbers are means of $\Snorm$ over the $3{,}896$ synthetic cards (lower is better; $1$ is M0), with $95\%$ quarter-block bootstrap
intervals where shown.

\subsection{The ablation ladder}

\begin{center}
{\small\begin{tabular}{lcccc}
\toprule
rung & all cards & $95\%$ CI & as-of $<2015$ & as-of $\ge2015$\\
\midrule
M0 replica & $1.0000$ & exact & $1.0000$ & $1.0000$\\
M1 driftless random walk (M0 width) & $1.0477$ & $[1.018,1.078]$ & $1.032$ & $1.067$\\
M2 variance term structure & $1.0643$ & $[1.036,1.094]$ & $1.058$ & $1.072$\\
honest belief of the shipped model ($\gamma=\lambda=0$) & $1.0453$ & $[1.018,1.075]$ & $1.033$ & $1.061$\\
\textbf{M11 full system (vintaged)} & $\mathbf{0.9680}$ & $\mathbf{[0.952,0.985]}$ & $0.9710$ & $\mathbf{0.9642}$\\
M11 + PIT calibration (M10) & $0.9707$ & $[0.956,0.991]$ & $0.9660$ & $0.9765$\\
\bottomrule
\end{tabular}}
\end{center}

Three facts stand out. First, the rungs that improve the \emph{honest} forecast (removing M0's noisy drift in M1, conditional volatility in M2)
make the ratio loss \emph{worse}: $1.048$ and $1.064$. Second, the full system's interval excludes $1$. Third, calibration is rejected
out of sample.

\paragraph{Anatomy of the honest-belief loss.} Breaking M1 down by template, the mean is inflated by a thin tail of clipped cards on single-asset,
long-horizon level cards (\code{ust\_single\_long} mean $1.224$, median $1.011$, $95$th percentile $3.05$). In each such case M0's drift, a
$300$-day trend extrapolated $126$--$189$ steps, happened to land near the outcome while the driftless centre did not. With comparable widths,
$\CRPS(\text{ours})/\CRPS(\MO)$ then approaches the clip. This is the convexity of Section~\ref{sec:ratio} made visible, and it is why the decision layer
exists. Retaining M0's drift does not fix it: with drift $\kappa=1$ the honest belief scores $1.035$, with $\kappa=0.5$ it scores $1.024$, because
the width then differs from M0's instead.

\subsection{The action grid}

On the same belief pools, scoring every $(\gamma,\lambda)$ (the floor not yet enforced):
\begin{center}
\begin{tabular}{lccl}
\toprule
belief & honest & best action & best score\\
\midrule
M1 (M0 width, driftless) & $1.0494$ & $(1.25,0.35)$ & $0.9554$\\
M2 (term structure, $\omega=0.5$) & $1.0641$ & $(1.50,0.35)$ & $0.9558$\\
drift $\kappa=0.5$, fat tails, residual bank & $1.0166$ & $(1.00,0.35)$ & $0.9692$\\
drift $\kappa=1$, fat tails, residual bank & $1.0244$ & $(1.50,0.50)$ & $0.9852$\\
\bottomrule
\end{tabular}
\end{center}
The decision layer is worth about $0.09$ on the same belief, the dominant effect in the system. Both parameters matter: $\gamma$ alone reaches about
$0.98$ at the grid edge, and $\lambda$ adds about $0.02$ by capping blow-ups (Theorem~\ref{thm:mix}).

\subsection{Tuning the belief}

One factor at a time around a base. The best fixed action is chosen on train (as-of and outcome before 2015) and the score is reported on test
(as-of 2015 or later). Round 1 did not enforce the spread floor; round 2 enforces it, with base $\omega=0$.
\begin{center}
{\footnotesize\setlength{\tabcolsep}{2.5pt}\begin{tabular}{lcccc|lccc}
\toprule
\multicolumn{5}{c|}{round 1 (no floor)} & \multicolumn{4}{c}{round 2 (floor $0.6$)}\\
variant & action & train & test & honest test & variant & action & train & test\\
\midrule
$\omega=0$ & $(1.5,0.2)$ & $\mathbf{0.9409}$ & $0.9618$ & $1.058$ & scale $1.1$ & $(1.0,0.1)$ & $\mathbf{0.9443}$ & $0.9658$\\
$\phi=0.995$ & $(1.5,0.35)$ & $0.9437$ & $0.9604$ & $1.068$ & \textbf{base} & $(1.25,0.2)$ & $0.9449$ & $0.9665$\\
$H=15$ & $(1.5,0.2)$ & $0.9440$ & $0.9671$ & $1.072$ & $\phi=0.995$ & $(1.0,0.2)$ & $0.9460$ & $0.9701$\\
no VR & $(1.25,0.2)$ & $0.9445$ & $0.9712$ & $1.074$ & $H=15$ & $(0.8,0.3)$ & $0.9465$ & $0.9674$\\
base $\omega=0.5$ & $(1.5,0.2)$ & $0.9449$ & $0.9675$ & $1.071$ & residual bank & $(1.0,0.3)$ & $0.9484$ & $0.9642$\\
scale $1.1$ & $(1.5,0.2)$ & $0.9452$ & $0.9599$ & $1.082$ & $\nu=12$ & $(1.0,0.3)$ & $0.9509$ & $0.9705$\\
$\nu=7$ & $(1.5,0.35)$ & $0.9508$ & $0.9738$ & $1.071$ &&&&\\
$\omega=1$ & $(1.5,0.35)$ & $0.9540$ & $0.9697$ & $1.089$ &&&&\\
\bottomrule
\end{tabular}}
\end{center}
The variants differ by less than the interval width. We keep the round-2 base: $\omega=0$ (long-run variance anchored at M0's window), EWMA
current variance, Gaussian, driftless, no residual bank. Choosing the action on test instead of train gives $0.962$ against $0.967$, so the action
choice is stable rather than overfit. Fat tails ($\nu=7$) are the worst belief variant: under the ratio loss, extra tail mass costs more CRPS on
ordinary cards than it saves in pinball on surprises.

\subsection{Out-of-sample performance by vintage}

\begin{center}
\begin{tabular}{lccccc}
\toprule
cards (as-of) & parameters & $n$ & mean & median & mean $(\gamma,\lambda)$ used\\
\midrule
2003--2008 & priors & $1083$ & $1.006$ & $0.959$ & $(0.50,0.20)$\\
2009--2014 & vintage $\le$2008 & $1075$ & $\mathbf{0.936}$ & $0.907$ & $(1.24,0.35)$\\
2015--2024 & vintage $\le$2014 + policy & $1738$ & $\mathbf{0.964}$ & $0.911$ & $(1.22,0.23)$\\
\bottomrule
\end{tabular}
\end{center}
$0.36\%$ of cards reach the clip. The priors row motivated raising the prior $\lambda$ to $0.4$: with nothing fitted, hedging harder is the
theory-consistent choice.

\subsection{By family on uniform dates}

\begin{center}
{\footnotesize\setlength{\tabcolsep}{3pt}\begin{tabular}{lcccc}
\toprule
family-shaped template & split & best action & best & at $(1,0.3)$\\
\midrule
F1 (long horizons) & $<2015$ / $\ge2015$ & $(1.25,0.4)$ / $(1.75,0.4)$ & $0.973$ / $0.983$ & $0.987$ / $1.017$\\
F2 ($63$--$64$ BD) & $<2015$ / $\ge2015$ & $(1.0,0.1)$ / $(1.25,0.3)$ & $0.910$ / $0.935$ & $0.917$ / $0.936$\\
F3 (joint) & $<2015$ / $\ge2015$ & $(1.25,0.2)$ / $(1.25,0.4)$ & $0.940$ / $0.960$ & $0.944$ / $0.962$\\
F4 ($21$ BD) & $<2015$ / $\ge2015$ & $(1.0,0.2)$ / $(1.0,0.2)$ & $0.953$ / $0.937$ & $0.954$ / $0.940$\\
\bottomrule
\end{tabular}}
\end{center}

\subsection{Stress windows (M12)}

\begin{center}
\begin{tabular}{lcccc}
\toprule
window & $n$ & M1 & M11 & M11 $90$th pct.\\
\midrule
GFC (2008-06 to 2009-06) & $208$ & $1.011$ & $1.013$ & $1.360$\\
taper (2013-04 to 2013-12) & $140$ & $1.012$ & $0.954$ & $1.142$\\
COVID (2020-01 to 2020-07) & $105$ & $1.049$ & $0.985$ & $1.438$\\
hikes (2022) & $177$ & $1.326$ & $1.062$ & $1.375$\\
worst decile of $|z_{\MO}|$ & $389$ & $0.970$ & $1.053$ & ---\\
\bottomrule
\end{tabular}
\end{center}
The system is weakest precisely where outcomes are far from M0's centre. Sharpening toward M0 costs most when M0 is badly wrong, and uniform-date
validation undersamples those cards.

\subsection{Event-selected practice cards}

On the practice units (realised values reconstructed from the published history; evaluation only):
\begin{center}
\begin{tabular}{lcccc}
\toprule
family & $n$ & best action (2024-vintage belief) & best & generic action $(1,0.3)$\\
\midrule
F1 & $18$ & $(1.25,0)$ & $0.871$ & $0.906$\\
F2 & $23$ & $(0.5,0.8)$ & $0.993$ & $1.035$\\
F3 & $20$ & $(0.5,0.3)$ & $0.937$ & $0.942$\\
F4 & $29$ & $(0,0.8)$ & $1.010$ & $1.085$\\
\bottomrule
\end{tabular}
\end{center}
F1's preference for $(1.25,0)$ contradicts the $394$ uniform-date F1-shaped cards, where that action scores $1.066$, so it was not adopted. For
F4, widening the belief is robust across actions: at scale $\times1.4$ the score is $0.968$--$0.973$ for every $\lambda$, and $0.973$ even for the
honest belief. For F2 the results are unstable (one combination reaches $0.922$ while its neighbours sit near $0.99$). These family blocks
(F4: $\times1.35$, action $(0,0.3)$; F2: $\times1.2$, action $(0.5,0.5)$; policy off) enter \emph{only} the 2024 vintage, because they were chosen
on outcomes up to 2024.

\paragraph{Text on practice cards (offline).} With the lexicon fallback only (no House endpoint off-platform), the numeric-only system scores
$1.0048$ and the text-on system $1.0038$ ($\Delta_{\text{text}}=+0.001$, helping on $17\%$ of cards). This is neutral, as designed for a capped fallback.

\subsection{The official Development submission}

The first Development upload (v1.0.1) was scored on $71$ cards, all admissible, with primary score $\mathbf{0.9782}$.
\begin{center}
\begin{tabular}{lcccc}
\toprule
family & $n$ & mean & median & share below $1$\\
\midrule
F1 & $18$ & $0.895$ & $0.881$ & $83\%$\\
F2 & $15$ & $0.965$ & $0.938$ & $67\%$\\
F3 & $16$ & $0.947$ & $0.893$ & $63\%$\\
F4 & $22$ & $1.078$ & $1.060$ & $32\%$\\
\bottomrule
\end{tabular}
\end{center}
\paragraph{Fidelity.} On the $65$ cards that the local evaluator can reconstruct, the official per-card scores and the local ones correlate at
$0.912$. The offline methodology tracks the official scorer.

\paragraph{The language model's contribution.} The official run called the House model and the local run could not, so the paired difference estimates
the LLM effect: $-0.0095$, with $95\%$ bootstrap interval $[-0.026,+0.009]$, better on $54\%$ of cards. This is consistent with a small benefit and
also with none. Text influence was therefore left unchanged.

\paragraph{The Final configuration on the same cards.} The Development cards are dated 2003--2024 and therefore use the 2008/2014 vintages, without the
F2/F4 blocks. Evaluated locally with the 2024 vintage and its family blocks (text off, hence conservative):
F1 $0.884$, F2 $0.957$, F3 $0.943$, F4 $1.006$, overall $\mathbf{0.952}$ against $0.983$ for the same cards under the Development configuration. The F2/F4
part of this estimate is in-sample and therefore optimistic.

\subsection{Robustness of the text path}

A fake House client returning adversarial answers on all $104$ practice units:
\begin{center}
{\small\begin{tabular}{p{0.34\textwidth}ccc}
\toprule
answer & admissible & path (full / numeric-only) & max sd/M0 sd (med., max)\\
\midrule
extreme (confidence $1$, extreme magnitude, both tails, $3$ scenarios, established) & $104/104$ & $102$ / $2$ & $1.71$, $3.55$\\
contradictory & $104/104$ & $102$ / $2$ & $1.11$, $3.55$\\
garbage (not JSON) & $104/104$ & $104$ / $0$ (lexicon) & $0.89$, $3.55$\\
\bottomrule
\end{tabular}}
\end{center}
The two numeric-only cases are the INR transfer cards, whose numeric belief is already about $3.5\times$ M0's scale. There the extreme text pushed
the spread beyond the $8\times$ sanity bound, and the update was rejected while the numeric forecast was kept.

\section{Engineering, safety and compliance}\label{sec:engineering}

\subsection{The fail-safe lattice}

The design goal is that no card can cost more than the baseline, whatever breaks. Let $\mathcal P$ denote the set of pipeline stages. Degradation
proceeds along a lattice of admissible outputs, from best to worst:
\begin{equation*}
  \text{full}\ \succ\ \text{numeric-only}\ \succ\ \text{M0 replica}\ \succ\ \text{minimal walk}\ \succ\ \text{constant around a stated anchor}.
\end{equation*}
\begin{itemize}[leftmargin=*]
\item \textbf{Text or fusion failure} (exception, timeout, unparseable answer, unreachable endpoint): $g=0$, and the result is P1 $=$ P0
(Propositions~\ref{prop:tilt} and~\ref{prop:mix}).
\item \textbf{Text update fails the spread sanity check} ($\sigma_c\notin[0.15,8]\cdot\sigma^{\MO}_c$ for some cell): the text update is dropped
and the numeric belief is kept.
\item \textbf{Any other exception}, including \code{KeyboardInterrupt} and a \code{SystemExit} raised deep in library code: the context is rebuilt
from the card and panels alone and M0's draws are written. This scores exactly $1$ (Proposition~\ref{prop:m0one}).
\item \textbf{M0 cannot be built}: a robust-scale walk with a common shock (perfect correlation is a safer joint default than independence for the
variogram when nothing is known).
\item \textbf{No panel at all}: small-spread draws around the level stated in the card's own prose (for example \code{CNY = 6.2097}), else $0$.
Admissible, and better than the $4$ of a missing card.
\end{itemize}

\subsection{Admissibility by construction}

The writer refuses to emit anything that would fail g0--g3: exactly three regular files; \code{draw} and \code{horizon} as \code{int32},
\code{value} as \code{float64}; every (draw, asset, horizon) exactly once, with draw identifiers $0,\dots,N-1$; finite values; positive spread in every cell;
$\max(200,n_{\min})\le N\le20000$; \code{unit\_id} and \code{asof} bound to the card; \code{representation} $=$ \code{samples}; a non-blank rationale.
Files are written under temporary names and atomically renamed, and temporaries are always removed, so a crash cannot leave a fourth file or a
half-written parquet.

\subsection{Runtime}

Pinned thread pools (four threads) respect the $256$-process limit. \code{HOME} and every cache point to \code{/tmp}, because the root filesystem
is read-only and user $65534$ has no home. There are no JIT caches (\code{/tmp} is non-executable). Measured on the practice units: $0.38$\,s mean
and $0.64$\,s maximum for the numeric path; about $2.7$\,s per unit in the container including start-up; House calls add at most $2\times120$\,s
under the per-call timeout. The phase clock of about $400$\,s per unit is therefore never binding.

\subsection{Determinism}

The run seed is $\mathrm{crc32}(\text{unit\_id})\oplus(\text{salt}\cdot2654435761)\oplus(97\cdot\text{QFBENCH\_SEED})$. Two runs with the same seed
produce identical parquet files (tested). The organizers' verification reruns on fresh seeds and compares bootstrap intervals, which the
Monte-Carlo sizes ($N\ge2000$, $N\ge4000$ for F4) make stable.

\subsection{Compliance matrix}

\begin{center}
\begin{tabular}{p{0.38\textwidth}p{0.55\textwidth}}
\toprule
rule & how it is met\\
\midrule
House model only, via the audited proxy & one endpoint from \code{MODEL\_ENDPOINT} with \code{MODEL\_TOKEN}; proxies read from the environment; no vendor APIs\\
Vendor tools off; pinned temperature and seed & no \code{tools} field; temperature $0$; seed pinned; thinking off\\
Request budget & at most $3$ calls against $25$; failures are not retried blindly (an admitted request costs a slot)\\
No neural checkpoints without approval & text is lexical; the foundation model is a no-op; no weights in the image\\
Fitted artifacts respect task cutoffs & vintaged with \code{fit\_cutoff} $<$ as-of (Section~\ref{sec:validation})\\
No cross-unit lookup & retrieval and analogues over the mounted unit only; no value derived from any unit's target is shipped\\
Text cutoff & corpus re-filtered at \code{min}(card, CLI) as-of; undated documents excluded\\
Joint draws & path-coherent generator; one scenario and one mixing variable per draw\\
Rationale & always written, including on fail-safe paths; states anchors, adjustments, documents cited, established versus inferred, and the decision layer\\
\bottomrule
\end{tabular}
\end{center}

\subsection{The audit (v1.0.1)}

A full audit after the first build found and fixed six issues. (1) A text update failing the sanity bound discarded the numeric forecast as
well; it now degrades to numeric-only. (2) Text widening is capped at $2\times$ and the maximum centre shift reduced ($\alpha$: $0.35\to0.25$).
(3) An HTTP $400$/$422$ now triggers one resend with the bare OpenAI body. (4) The last-resort constant is anchored on the level stated in the
card prose. (5) The rationale now reports $\lambda$. (6) Four unused configuration keys were removed, and a documentation count was corrected.
Earlier, adversarial tests had found two one-sidedness violations in schema coercion, a retrieval scoring flaw, a stale \code{default.yaml}
overriding code defaults (now guarded by a test that fails on any drift), and the \code{pandas}~3 datetime-resolution bug. The final suite has
$84$ tests (unit, integration including a mock House server, adversarial). Under the image's exact dependency stack (Python 3.13.5, numpy 2.1.3,
pandas 2.2.3, pyarrow 18.1.0) all pass, and a harness-exact container run (no network, read-only root, uid 65534, $64$\,MiB \code{/tmp}, $256$
processes) is admissible on $104/104$ units.

\section{Limitations and open questions}\label{sec:limitations}

\begin{enumerate}[leftmargin=*]
\item \textbf{The value of the language model is essentially unmeasured.} The House endpoint exists only inside the scoring environment. The
only estimate is the paired difference between the official run and a local text-off run: $-0.0095$ with interval $[-0.026,0.009]$. The gate
weights are therefore priors, chosen conservatively and protected by Propositions~\ref{prop:gate} and~\ref{prop:tilt}, not fitted.
\item \textbf{Event selection.} Uniform-date walk-forward undersamples the cards that the competition actually authors. The F2/F4 corrections come
from only $23$ and $29$ practice cards and are in-sample; the F4 widening is robust across actions, the F2 correction less so.
\item \textbf{The sharpening trade-off in stress.} The decision layer is weakest where M0 is badly wrong (worst decile $1.053$; 2022 hikes $1.062$).
A state-dependent action (less sharpening when the text flags stress) is the natural next step, but it can only be learned once text features
exist historically.
\item \textbf{A composite of three tilts.} Theorem~\ref{thm:bayes} gives the optimum for each component separately; the one-parameter tempered
operator is a compromise. A two- or three-parameter family (separate weights for the centre and the tails) might do better, at the cost of more
selection noise.
\item \textbf{Pinball non-convexity.} The mixing cap (Theorem~\ref{thm:mix}) is proved for CRPS and VS only.
\item \textbf{Transfer cards.} The EM premium ($1.25$) and the blend weight ($0.5$) are judgement calls; there is almost no emerging-market data to fit them.
\item \textbf{Cell ordering.} On three practice cards M0 enumerates cells in card order rather than sorted order; there the replica's draws differ
from M0's, though its law does not. The fail-safe then scores near, not exactly, $1$.
\end{enumerate}

\subsection{What would change the design}

A labelled corpus of dated text with outcomes \emph{before} each cutoff would allow fitting the gate (a logistic model for $\Pr(\Delta_{\text{text}}>0)$)
and adding a text action to the policy. Evidence that sealed cards are not event-selected would argue for the uniform-date actions. Evidence that
the pinball term dominates in the Final would argue for a separate tail tilt.

\appendix
\section{Further results and proofs}\label{app:proofs}

\subsection{Propriety of the scores}

\begin{proposition}[CRPS is strictly proper]\label{prop:crps-proper}
For distributions $F,G$ on $\R$ with finite first moments, with $\bar S(F,G)=\E_{Y\sim G}\CRPS(F,Y)$,
\begin{equation}
  \bar S(F,G)-\bar S(G,G)=\int_{-\infty}^{\infty}\bigl(F(t)-G(t)\bigr)^2\,dt\ \ge 0,
\end{equation}
with equality iff $F=G$.
\end{proposition}
\begin{proof}
Use the threshold representation $\CRPS(F,y)=\int(F(t)-\ind\{y\le t\})^2dt$. Taking $\E_G$ inside, for each $t$,
$\E_G(F(t)-\ind\{Y\le t\})^2=F(t)^2-2F(t)G(t)+G(t)=(F(t)-G(t))^2+G(t)(1-G(t))$. The last term does not depend on $F$ and integrates to
$\bar S(G,G)$. Equality forces $F(t)=G(t)$ for almost every $t$, hence everywhere by right-continuity.
\end{proof}

\begin{proposition}[The fair estimator is unbiased]
If $X_1,\dots,X_N$ are i.i.d.\ from $F$, then $\E$ of \eqref{eq:crps} equals $\CRPS(F,y)$ for every $y$.
\end{proposition}
\begin{proof}
$\E\frac1N\sum_k|X_k-y|=\E_F|X-y|$. The double sum over $k\ne l$ has $N(N-1)$ terms, each with expectation $\E_{F\otimes F}|X-X'|$, and the
diagonal terms are zero. Dividing by $N(N-1)$ is therefore unbiased, whereas $1/N^2$ would bias the spread term by the factor $(N-1)/N$.
\end{proof}

\begin{proposition}[Pinball is proper for the quantile]
For $Y\sim G$ continuous, $q\mapsto\E\rho_\alpha(Y-q)$ is convex with derivative $G(q)-\alpha$, so it is minimised at $q=G^{-1}(\alpha)$.
\end{proposition}
\begin{proof}
$\rho_\alpha(u)=\alpha u^+ +(1-\alpha)u^-$ is convex, so the expectation is convex in $q$. Differentiating under the integral,
$\frac{d}{dq}\E\rho_\alpha(Y-q)=-\alpha\PP(Y>q)+(1-\alpha)\PP(Y<q)=G(q)-\alpha$.
\end{proof}

\subsection{Why a ratio loss is not a difference loss}

\begin{proposition}[Lower expected CRPS does not imply a lower expected ratio]\label{prop:counter}
There exist $F$, $F^0$ and $G$ with $\E_G\CRPS(F,Y)<\E_G\CRPS(F^0,Y)$ but $\E_G[\CRPS(F,Y)/\CRPS(F^0,Y)]>1$.
\end{proposition}
\begin{proof}
Point masses have zero spread term, so $\CRPS(\delta_c,y)=|y-c|$. Let the baseline be $F^0=\delta_0$ and the forecast $F=\delta_1$, and let the
truth be the two-point law $G=0.8\,\delta_{0.45}+0.2\,\delta_{10}$. Then
\begin{align*}
  \E_G\CRPS(F,Y)&=0.8\cdot0.55+0.2\cdot9=2.24,\\
  \E_G\CRPS(F^0,Y)&=0.8\cdot0.45+0.2\cdot10=2.36,
\end{align*}
so $F$ is better in expected CRPS. But
\begin{equation*}
  \E_G\Bigl[\frac{\CRPS(F,Y)}{\CRPS(F^0,Y)}\Bigr]=0.8\cdot\frac{0.55}{0.45}+0.2\cdot\frac{9}{10}=1.158>1 .
\end{equation*}
\end{proof}
The mechanism is exactly the one seen in Section~\ref{sec:results}: when the baseline is sharp and right, its small denominator magnifies any
deviation, and outcomes near the baseline's centre receive large weight. This is what Theorem~\ref{thm:bayes} formalises.

\subsection{The clip}

\begin{lemma}
For any random variable $X\ge0$, $\E\min(4,X)\le\min(4,\E X)$, and the gap is at most $\E(X-4)^+$.
\end{lemma}
\begin{proof}
$x\mapsto\min(4,x)$ is concave, so Jensen gives the inequality. For the gap, $\min(4,\E X)-\E\min(4,X)\le\E X-\E\min(4,X)=\E(X-4)^+$.
\end{proof}
The clip bounds the damage from any single card and makes the loss even more asymmetric: the upside per card is at most about $0.7$, the downside at
most $3$.

\subsection{Resampling}

\begin{proposition}[Systematic resampling]
With normalised weights $p_1,\dots,p_M$ and $u\sim U[0,1)$, choose indices $i_k=\min\{i:\sum_{j\le i}p_j>(u+k)/N\}$, $k=0,\dots,N-1$. The number
$n_i$ of copies of $i$ satisfies $\E n_i=Np_i$ and $n_i\in\{\lfloor Np_i\rfloor,\lceil Np_i\rceil\}$.
\end{proposition}
\begin{proof}
The points $(u+k)/N$ form a grid of spacing $1/N$ with a uniform random offset. The interval $[\sum_{j<i}p_j,\sum_{j\le i}p_j)$ of length $p_i$
contains either $\lfloor Np_i\rfloor$ or $\lceil Np_i\rceil$ grid points, and in expectation exactly $Np_i$ of them, because the offset is uniform.
\end{proof}
Hence the resampled empirical law is an unbiased estimator of $Q_\gamma$ restricted to the pool, with the smallest possible variance per index,
compared with $Np_i(1-p_i)$ under multinomial resampling.

\begin{proposition}[Effective sample size]
If $p_i\propto\omega_i$, the variance of the self-normalised importance estimator of a bounded functional is, to first order, that of
$\mathrm{ESS}=(\sum p_i^2)^{-1}$ i.i.d.\ draws from the target.
\end{proposition}
\begin{proof}
Kong's approximation: the delta method applied to the ratio estimator gives
\begin{equation*}
  \Var\approx\frac{\Var_{Q}(f)}{M}\cdot\frac{\E_P[\omega^2]}{(\E_P\omega)^2},\qquad \frac{\E_P[\omega^2]}{(\E_P\omega)^2}\approx M\sum_i p_i^2 .
\end{equation*}
\end{proof}
The implementation rejects $\gamma$ when $\mathrm{ESS}<0.1M$.

\subsection{Variance reduction in the generator}

\begin{proposition}[Stratified scenarios]
Let the Monte-Carlo estimator of $\E f(\bm Y)$ under a mixture $\sum_j\pi_jP_j$ allocate exactly $N\pi_j$ draws to component $j$. Its variance is
$\frac1N\sum_j\pi_j\Var_{P_j}(f)$, never larger than the multinomial-allocation variance $\frac1N\Var_{\sum\pi_jP_j}(f)$.
\end{proposition}
\begin{proof}
Law of total variance: $\Var_{\text{mix}}(f)=\sum_j\pi_j\Var_{P_j}(f)+\sum_j\pi_j(\E_{P_j}f-\bar f)^2$. Stratification removes the second, between-component term.
\end{proof}

\begin{proposition}[Antithetic pairs]
Pairing each increment vector $\bm\xi$ with $-\bm\xi$ under a shared scenario and mixing variable makes the sample mean of the base layer equal the
centre exactly, and for any odd functional the Monte-Carlo error is zero.
\end{proposition}
\begin{proof}
The base layer is $\mu+\bm b+\bm c\sqrt W X(\bm\xi)$ with $X$ linear in $\bm\xi$. Summing over a pair cancels the linear term, and an odd functional
of the centred layer cancels pairwise.
\end{proof}

\subsection{EWMA volatility}

\begin{lemma}
The recursion $\hat\sigma^2_t=\lambda\hat\sigma^2_{t-1}+(1-\lambda)x_t^2$ is the weighted average $\sum_{i\ge0}(1-\lambda)\lambda^ix_{t-i}^2$, with
weights summing to one and centre of mass $\lambda/(1-\lambda)$ steps in the past. With $\lambda=2^{-1/H}$ the weight halves every $H$ steps; for
$H=30$, $\lambda=0.9772$ and the effective window is about $43$ steps.
\end{lemma}
\begin{proof}
Unroll the recursion. The mean lag is $\sum_i i(1-\lambda)\lambda^i=\lambda/(1-\lambda)$, which for $\lambda=0.9772$ is about $42.8$.
\end{proof}
If $x_t$ has constant variance $\sigma^2$, the estimator is unbiased, and its variance is $\frac{1-\lambda}{1+\lambda}\Var(x^2)$, about
$0.0115\,\Var(x^2)$ for $H=30$: roughly the precision of $87$ independent squared steps, a sensible compromise between reactivity and noise.

\subsection{The variance ratio}

\begin{proposition}
For a covariance-stationary step process with autocorrelations $\rho_k$,
\begin{equation}
  VR(s)=\frac{\Var(\sum_{i=1}^sx_i)}{s\Var(x)}=1+2\sum_{k=1}^{s-1}\Bigl(1-\frac ks\Bigr)\rho_k .
\end{equation}
For AR(1) steps with $\rho_k=\phi^k$, $VR(s)\to\frac{1+\phi}{1-\phi}$ as $s\to\infty$.
\end{proposition}
\begin{proof}
$\Var(\sum x_i)=\sum_{i,j}\Cov(x_i,x_j)=s\Var(x)+2\sum_{k=1}^{s-1}(s-k)\rho_k\Var(x)$; divide by $s\Var(x)$. For AR(1) the Ces\`aro-weighted sum
tends to $\sum_{k\ge1}\phi^k=\phi/(1-\phi)$, giving $1+2\phi/(1-\phi)$.
\end{proof}
So $VR>1$ indicates trending steps (momentum) and $VR<1$ mean reversion at the horizon. Since the estimator's standard error scales like
$\sqrt{s/n}$, the pseudo-count shrinkage \eqref{eq:vr} is the appropriate response to its noise.

\subsection{Isotonic regression}

\begin{proposition}[PAVA]
The pool-adjacent-violators algorithm returns the Euclidean projection of $\bm y\in\R^n$ onto the monotone cone
$\{\bm z:z_1\le\dots\le z_n\}$.
\end{proposition}
\begin{proof}
Pooling a violating adjacent pair into its weighted mean never increases the squared distance to any feasible point that is constant on the
pooled block. The algorithm terminates at a feasible point satisfying the KKT conditions of the convex quadratic program: blocks are level sets,
and each block value is the mean of its data.
\end{proof}

\subsection{Termination and safety of the decision layer}

\begin{proposition}
For a fixed action $(\gamma,\lambda)$ the stepping rule terminates after at most $\lceil4\gamma\rceil+1$ candidates and always returns admissible
draws.
\end{proposition}
\begin{proof}
Each rejection lowers $\gamma$ by $0.25$, and $\gamma=0$ is accepted unconditionally: it is a plain resample of the honest belief, which already
passed the spread sanity check. Mixing with M0 draws preserves finiteness and positive spread.
\end{proof}

\section{Parameters}\label{app:params}

Built-in defaults are \emph{priors}: no value in this table was fitted, unless the table says it was. Fitted values live in the vintaged artifacts.

\begin{center}
\small
\begin{longtable}{p{0.33\textwidth}p{0.17\textwidth}p{0.42\textwidth}}
\toprule
parameter & value & role / provenance\\
\midrule
\endhead
M0 window, draws, override ratio & $300$, $500$, $2$ & official specification\\
EWMA half-life $H$ & $30$ BD & prior\\
long window $L$ & $1260$ BD & prior\\
long-run blend $\omega$ & $0.5$ prior; $0$ fitted & vintages $\le$2014, $\le$2024 (tuning round 1)\\
persistence $\phi$ & $0.985$ / BD & prior (half-life about $46$ BD)\\
global scale $c$ & $1.0$ & prior; within noise in tuning\\
VR pseudo-count, clip & $40$; $[0.6,1.6]$ & prior\\
drift shrink (level / return) & $0$ / $0.25$ prior; $0$ / $0$ fitted & ratio convexity\\
scale-mixture dof $\nu$ & $12$ prior; $\approx\infty$ fitted & fat tails rejected\\
correlation blend $\eta$ & $0.15$ & prior\\
soft floor $f_0$, $k$ & $-0.10$, $0.25$ & UST only\\
macro trend half-life, weight & $6$ months, $0.7$ & prior\\
transfer blend, EM premium & $0.5$, $1.25$ & judgement\\
regime thresholds, slope & $0.8$, $1.35$, $k=8$ & prior\\
residual bank & off & no robust gain\\
passage size, top-$k$, budget & $900$, $12$, $14{,}000$ chars & engineering\\
dedup Jaccard, window & $0.55$, $5$ days & prior\\
BM25 $k_1$, $b$; recency half-life & $1.4$, $0.75$; $45$ days & standard\\
retrieval weights (rec, asset, nov, type) & $0.6$, $0.8$, $0.3$, $0.4$ & prior\\
LLM calls, retries, max tokens, temperature & $3$, $1$, $1800$, $0$ & budget discipline\\
gate $(w_0,\dots,w_4)$ & $-1.2$, $1.1$, $2.2$, $0.35$, $1.4$ & prior (no labelled text)\\
gate cap, lexicon cap & $0.85$, $0.25$ & prior\\
family $(\delta_f,\theta_f)$ & F1 $(-0.3,0.8)$, F2 $(0.4,1)$, F3 $(0,0.8)$, F4 $(0.2,0.9)$ & prior\\
update $\alpha,\beta,\lambda_\kappa,\kappa^\star$ & $0.25$, $0.30$, $0.15$, $0.55$ & prior; $\alpha$ reduced in audit\\
text vol cap & $2\times$ & audit\\
tail weight max, tail scale & $0.12$, $2.2$ & prior\\
decision $(\gamma,\lambda)$ & prior $(0.5,0.4)$; $\le$2008 $(1.25,0.35)$; $\le$2014 $(1.25,0.2)$; $\le$2024 $(1.0,0.3)$ & fitted per vintage\\
F4 / F2 blocks ($\le$2024 only) & $\times1.35$, $(0,0.3)$ / $\times1.2$, $(0.5,0.5)$ & practice cards (in-sample)\\
spread floor, ESS floor & $0.6$, $0.1M$ & honesty constraints\\
policy ridge, pessimism & $100$, $2$ & fitted per vintage; OOS evidence\\
belief pool, evaluation outcomes, $N$ & $6000$, $3000$, $2000$ ($4000$ F4) & Monte-Carlo accuracy\\
\bottomrule
\end{longtable}
\end{center}

\section{Reproduction}\label{app:repro}

All commands run from the repository root \path|agenthon-track2|. The practice units live in the organizers' public repository.

\begin{enumerate}[leftmargin=*]
\item Rebuild history: \path|python offline/data/build_history.py --units <units> --out offline/data|.
\item Generate cards: \path|python offline/walk_forward/generate_cutoffs.py| ($3{,}896$ cards).
\item Sanity rung: \path|python -m offline.ablations.M0_baseline 8| (must print $1.0000$).
\item Ladder: \path|python -m offline.ablations.M1_random_walk 8|, \path|... M2_statistical 8|.
\item Tuning: \path|python offline/experiments/tune.py --prefix R2_ --base '{"numerical":{"lr_omega":0.0}}'|.
\item Policy: \path|python offline/rl/train_policy.py --run R2_base|.
\item Production-exact walk-forward: \path|python offline/walk_forward/run.py --name M11_full --overrides '{"_vintaged": true}'|.
\item Calibration ablation: \path|python offline/experiments/fit_calibration.py --run M11_full|.
\item Practice evaluation: \path|python offline/walk_forward/practice_eval.py --rung M11_full|;
      family actions: \path|python offline/walk_forward/practice_actions.py|.
\item Tables: \path|python offline/walk_forward/aggregate.py M0_baseline M1_random_walk M11_full --save ablation_ladder|;
      stress: \path|python offline/walk_forward/stress.py M11_full|.
\item Tests: \path|python -m pytest| ($84$ tests). Official gates on all units: \path|python scripts/run_local.py --units <units> --no-llm| then
      \path|python scripts/validate_submission.py --units <units> --outputs artifacts/forecasts|.
\item Container: \path|docker build --platform linux/amd64 -t agenthon-t2:1.0.1 .| and
      \path|powershell -File scripts/docker_smoke.ps1 -All| ($104/104$).
\item Adversarial LLM stress: \path|python scripts/llm_stress.py --units <units> --mode extreme|.
\end{enumerate}

\paragraph{Submitted artifact.} Image \path|docker.io/tps3108/agenthon-t2|, linux/amd64 digest
\path|sha256:00e0ee761cdbabca422156a3bfb096bdddd35783d3fb98dfb09fbafa2a0b98bb|, disclosed House model
\path|nvidia/nemotron-3-super-120b-a12b| (\path|rl-030326-fp8|, training cutoff unpublished).

\section{Further theory}\label{app:theory}

\subsection{Closed-form expected CRPS for Gaussian beliefs}

\begin{lemma}\label{lem:absnormal}
If $U\sim\mathcal N(m,v)$ then $\E|U|=\sqrt v\,\bigl[2\varphi(m/\sqrt v)+(m/\sqrt v)\bigl(2\Phi(m/\sqrt v)-1\bigr)\bigr]$.
\end{lemma}
\begin{proof}
Write $U=m+\sqrt vZ$ and split at $Z=-m/\sqrt v$. We have $\E(m+\sqrt vZ)^+=m\Phi(m/\sqrt v)+\sqrt v\varphi(m/\sqrt v)$ and
$\E(m+\sqrt vZ)^-=-m\Phi(-m/\sqrt v)+\sqrt v\varphi(m/\sqrt v)$; adding them gives the result.
\end{proof}

\begin{proposition}[Expected CRPS of a Gaussian forecast under a Gaussian truth]\label{prop:ecrps}
Let the truth be $Y\sim\mathcal N(0,1)$ and the forecast $F=\mathcal N(m,\sigma^2)$. Then
\begin{equation}
  \E_Y\CRPS(F,Y)=\sqrt{1+\sigma^2}\;\Bigl[2\varphi(\tilde m)+\tilde m\bigl(2\Phi(\tilde m)-1\bigr)\Bigr]-\frac{\sigma}{\sqrt\pi},
  \qquad \tilde m=\frac{m}{\sqrt{1+\sigma^2}} .
\end{equation}
It is minimised at $(m,\sigma)=(0,1)$, where $\E\CRPS=1/\sqrt\pi\approx0.5642$.
\end{proposition}
\begin{proof}
$\E|X-Y|$ with $X\perp Y$ is $\E|U|$ for $U=X-Y\sim\mathcal N(m,1+\sigma^2)$, given by Lemma~\ref{lem:absnormal}. $\E|X-X'|$ for independent copies is
$\E|\mathcal N(0,2\sigma^2)|=2\sigma/\sqrt\pi$, and half of it is subtracted. At $m=0$, $\sigma=1$: $\sqrt2\cdot2\varphi(0)-1/\sqrt\pi=2/\sqrt\pi-1/\sqrt\pi=1/\sqrt\pi$.
Propriety (Proposition~\ref{prop:crps-proper}) gives the minimality.
\end{proof}
Around the optimum the loss is quadratic in the centre error. With $\sigma=1$ let $f(m)=\E\CRPS$ and $g(u)=2\varphi(u)+u(2\Phi(u)-1)$, so
$f(m)=\sqrt2\,g(m/\sqrt2)-1/\sqrt\pi$. Since $g'(u)=2\Phi(u)-1$ and $g''(u)=2\varphi(u)$, we get $f'(0)=0$ and $f''(0)=\sqrt2\cdot\frac12\cdot2\varphi(0)=1/\sqrt\pi$, hence
\begin{equation}
  \E\CRPS\;\approx\;\frac{1}{\sqrt\pi}\Bigl(1+\frac{m^2}{2}\Bigr).
\end{equation}
A centre error of $0.65$ horizon standard deviations, the size of M0's drift noise at $126$ steps (Section~\ref{sec:numerical}), costs about $21\%$ of
expected CRPS; the exact value from Proposition~\ref{prop:ecrps} is $0.6813$ against $0.5642$, that is $+20.8\%$. Removing it therefore helps expected CRPS, as the ladder confirms
for raw scores. The next result explains why it nonetheless does not help the \emph{ratio}.

\subsection{A second-order expansion of the expected ratio}

\begin{theorem}[Delta method for the per-card ratio]\label{thm:ratioexp}
Let $S=S(Q,\bm Y)$ and $S^0=S(F^0,\bm Y)$ be positive with means $\mu,\mu_0$, and let $c_0=\sd(S^0)/\mu_0$ be the baseline's coefficient of
variation across outcomes and $\rho=\mathrm{Corr}(S,S^0)$. To second order,
\begin{equation}\label{eq:ratioexp}
  \E\Bigl[\frac{S}{S^0}\Bigr]\;\approx\;\frac{\mu}{\mu_0}\Bigl(1+c_0^2-\rho\,c_0\,c\Bigr),\qquad c=\sd(S)/\mu .
\end{equation}
\end{theorem}
\begin{proof}
Expand $g(s,s_0)=s/s_0$ around $(\mu,\mu_0)$: $g_{s}=1/\mu_0$, $g_{s_0}=-\mu/\mu_0^2$, $g_{ss}=0$, $g_{s s_0}=-1/\mu_0^2$, $g_{s_0s_0}=2\mu/\mu_0^3$.
Taking expectations kills the first-order terms, and
$\E g\approx\frac{\mu}{\mu_0}+\frac12\bigl(2\cdot(-\tfrac1{\mu_0^2})\Cov(S,S^0)+\tfrac{2\mu}{\mu_0^3}\Var S^0\bigr)
=\frac{\mu}{\mu_0}\bigl(1-\frac{\Cov(S,S^0)}{\mu\mu_0}+\frac{\Var S^0}{\mu_0^2}\bigr)$. Substitute $\Cov=\rho\,\sd(S)\sd(S^0)$.
\end{proof}

\begin{corollary}[Why hedging toward the baseline pays]
For a fixed ratio of means $\mu/\mu_0$, the expected ratio decreases in the correlation $\rho$ between our score and the baseline's across
outcomes. $\rho=1$ when $Q=F^0$, and it falls as the centre or scale of $Q$ moves away from $F^0$. Moving the centre away from M0's therefore has
two effects: it lowers $\mu$ (good, if the move is right on average) and it lowers $\rho$ (bad). With $c_0$ of order one, typical for CRPS across
outcomes, the second effect is first-order in size. This is the analytic content of the observation that the honest, driftless belief scores
$1.048$ although it is better in raw CRPS.
\end{corollary}
\begin{remark}
For a Gaussian baseline and truth, $S^0=\CRPS(F^0,Y)$ has a coefficient of variation of $0.715$ (computed by simulation with $2\times10^6$ outcomes), so $c_0^2\approx0.51$. Values of $c_0^2$ near $0.5$ mean that
denominators near their minimum ($0.234\sigma$, Proposition~\ref{prop:shape}) contribute disproportionately, and the tilt of
Theorem~\ref{thm:bayes} is exactly the reweighting that corrects for this.
\end{remark}

\subsection{The Gaussian variogram in closed form}

\begin{proposition}\label{prop:gvs}
If $(X_i,X_j)$ is bivariate Gaussian with $\Var(X_i-X_j)=\sigma^2_{ij}$ and $\E(X_i-X_j)=\delta_{ij}$, then for $\delta_{ij}=0$,
\begin{equation}
  \E|X_i-X_j|^p=\sigma_{ij}^p\,\frac{2^{p/2}\,\Gamma\bigl(\frac{p+1}{2}\bigr)}{\sqrt\pi},\qquad
  p=\tfrac12:\ \E|X_i-X_j|^{1/2}=\sigma_{ij}^{1/2}\cdot\frac{2^{1/4}\Gamma(3/4)}{\sqrt\pi}\approx0.8222\,\sigma_{ij}^{1/2}.
\end{equation}
\end{proposition}
\begin{proof}
$X_i-X_j=\sigma_{ij}Z$, and $\E|Z|^p=\frac{2^{p/2}\Gamma((p+1)/2)}{\sqrt\pi}$ from the half-normal moment integral
$2\int_0^\infty z^p\varphi(z)dz$ with the substitution $u=z^2/2$. Numerically, $2^{1/4}\cdot1.2254/1.7725=0.8222$.
\end{proof}

\begin{corollary}[The cost of independent sampling]
With $\sigma_{ij}^2=\sigma_i^2+\sigma_j^2-2\rho_{ij}\sigma_i\sigma_j$, sampling the cells independently replaces $\sigma_{ij}$ by
$\sqrt{\sigma_i^2+\sigma_j^2}$. For $\sigma_i=\sigma_j$ and $\rho=0.9$ (neighbouring Treasury tenors) the predicted $\E|X_i-X_j|^{1/2}$ is inflated by
$(2/0.2)^{1/4}=1.78$. By Proposition~\ref{prop:vs-opt} each such pair then contributes a squared bias of order $(0.78\cdot0.822)^2\sigma_{ij}$ to the
expected variogram score. For one asset at horizons $s_1<s_2$ under the path structure, $\sigma^2_{ij}=V(s_2)-V(s_1)$, while independent draws give
$V(s_1)+V(s_2)$: an inflation factor of $\bigl((V_1+V_2)/(V_2-V_1)\bigr)^{1/4}$, which is $2^{1/4}\approx1.19$ for $s=(21,63)$ under a random walk.
\end{corollary}

\subsection{Monte-Carlo error in the tails}

\begin{proposition}[Asymptotic variance of an empirical quantile]
For i.i.d.\ draws from a continuous $F$ with density $f$ positive at $q_\alpha$,
$\sqrt N(\hat q_\alpha-q_\alpha)\Rightarrow\mathcal N\bigl(0,\alpha(1-\alpha)/f(q_\alpha)^2\bigr)$.
\end{proposition}
\begin{proof}
Bahadur representation: $\hat q_\alpha-q_\alpha=-\frac{\hat F(q_\alpha)-\alpha}{f(q_\alpha)}+o_P(N^{-1/2})$, and $N\hat F(q_\alpha)\sim\mathrm{Bin}(N,\alpha)$.
\end{proof}
For a Gaussian cell at $\alpha=0.01$, $f(q_\alpha)=\varphi(2.326)=0.0267$, so $\sd(\hat q_{0.01})\approx\sqrt{0.0099/N}/0.0267$ standard deviations:
$0.264$ at $N=200$, $0.167$ at $N=500$, $0.083$ at $N=2000$ and $0.059$ at $N=4000$. Because the pinball loss is locally linear in the quantile
error with slope bounded by $\max(\alpha,1-\alpha)$, this Monte-Carlo noise enters the tail component directly. F4 therefore uses
$N\ge4000$ and all other cards $N\ge2000$, well above the admissibility floor of $200$.

\begin{proposition}[Monte-Carlo variance of the CRPS estimator]
For fixed $y$ the estimator \eqref{eq:crps} is a U-statistic of degree two with variance $O(1/N)$. Its leading term is
$\frac{1}{N}\Var_F\bigl(|X-y|-\E_{X'}|X-X'|\bigr)$.
\end{proposition}
\begin{proof}
Hoeffding decomposition: the first-order projection of the kernel $h(x,x')=\frac12(|x-y|+|x'-y|)-\frac12|x-x'|$ is
$h_1(x)=|x-y|/2+\E|X-y|/2-\E_{X'}|x-X'|/2$, and $\Var=\frac{4}{N}\Var h_1+O(N^{-2})$.
\end{proof}

\subsection{GARCH correspondence}

\begin{proposition}
For GARCH(1,1), $\sigma^2_{t+1}=\omega+\alpha x_t^2+\beta\sigma^2_t$ with $\alpha+\beta<1$, the forecast satisfies
$\E_t\sigma^2_{t+i}=\bar\sigma^2+(\alpha+\beta)^{i-1}(\sigma^2_{t+1}-\bar\sigma^2)$ with $\bar\sigma^2=\omega/(1-\alpha-\beta)$.
\end{proposition}
\begin{proof}
$\E_t x^2_{t+i-1}=\E_t\sigma^2_{t+i-1}$, so $\E_t\sigma^2_{t+i}=\omega+(\alpha+\beta)\E_t\sigma^2_{t+i-1}$, a first-order linear recursion whose fixed
point is $\bar\sigma^2$.
\end{proof}
Proposition~\ref{prop:agg} is therefore the aggregated GARCH forecast with $\phi=\alpha+\beta$, $\sigma^2_{\text{now}}$ playing the role of
$\sigma^2_{t+1}$ (estimated by EWMA, the IGARCH limit) and $\sigma^2_{\mathrm{LR}}$ that of $\bar\sigma^2$ (estimated by window variances rather
than by maximum likelihood). Replacing likelihood estimation with this robust plug-in avoids the convergence failures of per-card GARCH fits on
short or irregular panels, at no cost to the fail-safe.

\subsection{Scenarios and the variogram}

\begin{proposition}
Let a mixture have components that differ only by a \emph{common} shift $b$ of all cells. Then the variogram statistic
$\E|X_i-X_j|^p$ equals that of the unshifted component. If the shift differs across cells, the pairwise statistic becomes a mixture of shifted
moments, which Proposition~\ref{prop:gvs} generalises through the folded-normal moment with non-zero mean $\delta_{ij}$.
\end{proposition}
\begin{proof}
$X_i+b-(X_j+b)=X_i-X_j$. With cell-specific shifts the difference in component $j$ is $\mathcal N(\delta_{ij},\sigma_{ij}^2)$, whose $p$-th absolute
moment increases in $|\delta_{ij}|$.
\end{proof}
This is why text scenarios are applied as coherent joint stories: a "hawkish surprise" scenario moving all tenors in the same direction leaves the
cross-tenor geometry intact and only moves the marginals, while a steepening story deliberately changes it.


\end{document}
